\chapter{PHƯƠNG PHÁP ĐỀ XUẤT}
\label{ch:phuong_phap_de_xuat}

\vspace{-1.5cm}

\begin{quotation}
    \noindent
    \small\textit{
        Chương này trình bày phương pháp được sử dụng trong khóa luận. Trọng tâm là kế thừa AutoShot, giữ lại mạng xương sống đã được NAS tối ưu, sau đó cải tiến tầng phân loại, hàm mất mát, nhãn giám sát và các bước hậu xử lý xác suất để hình thành AutoShotV2. Các phần nền tảng về NAS và TransNetV2 được giữ lại ở mức cần thiết nhằm làm rõ cơ sở kế thừa trước khi đi vào phương pháp cải tiến; dữ liệu, thiết lập thực nghiệm và kết quả định lượng được trình bày ở Chương~\ref{ch:bo_du_lieu_va_thuc_nghiem}.
    }
\end{quotation}

\vspace{1em}

\newcolumntype{Q}{>{\raggedright\arraybackslash}X}

\section{Tổng quan hướng tiếp cận}
Khóa luận tiếp cận bài toán theo hướng kế thừa một mạng xương sống (backbone) phát hiện ranh giới cảnh quay (SBD) đã được tối ưu bằng tìm kiếm kiến trúc mạng (Neural Architecture Search - NAS), sau đó tập trung cải thiện các thành phần ở tầng đầu ra. Cụ thể, AutoShot được dùng làm mô hình nền tảng; phần mạng xương sống được giữ nguyên để bảo toàn lợi ích của quá trình tìm kiếm kiến trúc, còn AutoShotV2 bổ sung tầng phân loại (ClassificationHead) mới, hàm mất mát Focal (Focal Loss), nhãn nhiều điểm dương (many-hot), hiệu chỉnh nhiệt độ (temperature scaling) và làm mượt Gaussian (Gaussian smoothing) nhằm xử lý mất cân bằng lớp, tăng tín hiệu giám sát quanh ranh giới và ổn định chuỗi xác suất dự đoán theo thời gian.

\noindent
Hai mục tiếp theo hệ thống hóa hai khối kiến thức nền tảng được khóa luận kế thừa: \textbf{(i)} các nguyên lý và kỹ thuật của Neural Architecture Search (Mục~\ref{sec:nas}) -- là công cụ trung tâm mà AutoShot vận dụng để tìm ra kiến trúc tối ưu cho SBD, và \textbf{(ii)} mô hình TransNetV2 (Mục~\ref{sec:transnetv2}) -- là mô hình cơ sở trực tiếp để khóa luận so sánh và cải tiến. Sơ đồ luồng xử lý và kiến trúc tổng quan của toàn bộ hệ thống được minh họa chi tiết trong Hình~\ref{fig:overall_flowchart}.

\begin{figure}[H]
    \centering
    \includegraphics[width=0.92\textwidth]{images/Flowchart.png}
    \caption[Sơ đồ luồng xử lý tổng quan của hệ thống]{Sơ đồ luồng xử lý tổng quan của hệ thống phát hiện ranh giới cảnh quay AutoShot/AutoShotV2. Video đầu vào trải qua bước tiền xử lý, trích xuất đặc trưng kết hợp với nhánh tính độ tương đồng để tạo vector đặc trưng tổng hợp $4864$-chiều, sau đó qua các lớp kết nối đầy đủ và phân loại song song để dự đoán xác suất ranh giới.}
    \label{fig:overall_flowchart}
\end{figure}

\section{Kỹ thuật tìm kiếm kiến trúc mạng nơ-ron}
\label{sec:nas}

Tìm kiếm kiến trúc mạng nơ-ron (Neural Architecture Search — NAS) là quá trình tự động hóa việc thiết kế kiến trúc mạng nơ-ron nhằm đạt hiệu suất cao trên một nhiệm vụ cụ thể mà không cần can thiệp thủ công từ chuyên gia \cite{white2023_nas1000}. NAS được xem là bước tiến tự nhiên của học máy tự động (Automated Machine Learning — AutoML), mở rộng phạm vi tối ưu hóa từ siêu tham số sang toàn bộ cấu trúc mô hình. Trong thập kỷ qua, các kiến trúc được tìm kiếm tự động liên tục thiết lập kết quả tốt nhất trên nhiều bộ chuẩn đánh giá, đáng chú ý nhất là bộ dữ liệu ImageNet. Lĩnh vực này phát triển nhanh chóng với hơn 1000 bài báo được công bố chỉ trong hai năm tính đến 2022 \cite{white2023_nas1000}, phản ánh tầm quan trọng của tự động hóa thiết kế mô hình trong học sâu hiện đại.

Theo khảo sát toàn diện của Elsken et al.\ \cite{elsken2018_nas_survey} và Ren et al.\ \cite{ren2020_nas_survey}, hầu hết các phương pháp NAS đều được tổ chức quanh ba thành phần cốt lõi liên kết chặt chẽ với nhau: \textit{không gian tìm kiếm} xác định phạm vi các kiến trúc có thể xét đến, \textit{chiến lược tìm kiếm} quyết định cách duyệt qua không gian đó một cách hiệu quả, và \textit{chiến lược ước tính hiệu suất} giúp đánh giá nhanh chất lượng từng ứng viên mà không phải huấn luyện đầy đủ từ đầu. Ba thành phần này tạo thành một vòng lặp phản hồi: thiết kế không gian tìm kiếm ảnh hưởng trực tiếp đến chiến lược tìm kiếm phù hợp, còn phương pháp ước tính hiệu suất quyết định tốc độ toàn bộ vòng lặp NAS. Hình~\ref{fig:nas_pipeline} minh họa sơ đồ tổng quan của ba thành phần này.

\begin{figure}[H]
    \centering
    \begin{tikzpicture}[>=Latex,font=\footnotesize]
        \node[draw, rounded corners, fill=blue!12, text width=3.6cm, minimum height=1.25cm, align=center] (space) at (0,0) {Không gian tìm kiếm\\(Search Space)\\$\mathcal{A}$};
        \node[draw, rounded corners, fill=green!12, text width=3.6cm, minimum height=1.25cm, align=center] (search) at (5.0,0) {Chiến lược tìm kiếm\\(Search Strategy)};
        \node[draw, rounded corners, fill=orange!18, text width=3.9cm, minimum height=1.25cm, align=center] (est) at (2.5,-2.8) {Ước tính hiệu suất\\(Performance Estimation)};
        \node[draw, rounded corners, fill=purple!12, text width=3.0cm, minimum height=1.25cm, align=center] (output) at (10.2,-1.4) {Kiến trúc tối ưu\\$a^*$};

        \draw[->, thick] (space) -- node[above] {ứng viên} (search);
        \draw[->, thick] (search.south west) -- node[right] {đánh giá} (est.north east);
        \draw[->, thick] (est.north west) -- node[left] {phản hồi} (space.south east);
        \draw[->, thick] (search.east) -- node[above] {chọn tốt nhất} (output.north west);
        \draw[dashed, thick] (output.south) -- ++(0,-0.9) node[below] {ngân sách $t$};
    \end{tikzpicture}
    \caption[Sơ đồ ba thành phần cốt lõi của NAS]{Sơ đồ ba thành phần cốt lõi của NAS: không gian tìm kiếm (không gian tìm kiếm), chiến lược tìm kiếm (search strategy) và chiến lược ước tính hiệu suất (performance estimation strategy) liên kết nhau thành vòng lặp phản hồi. Đầu ra là kiến trúc tối ưu $a^*$ đáp ứng ngân sách tính toán $t$.}
    \label{fig:nas_pipeline}
\end{figure}

\subsection{Định nghĩa bài toán NAS}
\label{subsec:nas_def}

NAS tìm kiến trúc $a^*$ tối ưu hóa độ chính xác trên tập xác thực trong ngân sách tính toán $t$:
\begin{equation}
    a^* = \arg\max_{a \in \mathcal{A}} \; \mathrm{val\_acc}(a), \quad \text{trong ngân sách } t.
    \label{eq:nas_objective}
\end{equation}
Thách thức: không gian $\mathcal{A}$ rời rạc, không khả vi, có thể hàng triệu cấu hình — huấn luyện đầy đủ từng kiến trúc không khả thi. Trong SBD, kiến trúc tối ưu phụ thuộc mạnh vào đặc điểm dữ liệu (nhịp cắt, loại chuyển cảnh), nên NAS đặc biệt có giá trị để tìm mạng xương sống chuyên biệt cho từng phân phối dữ liệu — đây là luận điểm trung tâm của AutoShot.

\subsection{Không gian tìm kiếm}
\label{subsec:search_space}

Không gian tìm kiếm $\mathcal{A}$ xác định tập hợp tất cả các kiến trúc mạng nơ-ron khả thi. Thiết kế không gian tìm kiếm là quyết định quan trọng nhất trong NAS vì nó định hình toàn bộ quá trình: không gian quá lớn dẫn đến chi phí tìm kiếm khổng lồ và khó hội tụ, trong khi không gian quá hẹp có thể loại bỏ các kiến trúc tiềm năng. Thực tế, phần lớn các kiến trúc tìm kiếm tốt đều được thiết kế không gian tìm kiếm lấy cảm hứng mạnh từ các kiến trúc thủ công tiên tiến \cite{white2023_nas1000}.

Ba dạng không gian tìm kiếm phổ biến: (i) \textbf{Macro} — tìm kiếm toàn bộ mạng, tự do tối đa nhưng chi phí rất cao; (ii) \textbf{Hierarchical} — xây dựng kiến trúc theo nhiều cấp lồng nhau, cân bằng giữa linh hoạt và chi phí; (iii) \textbf{Cell-based} — tìm một đơn vị cấu trúc nhỏ (cell/block) rồi xếp chồng, phổ biến nhất hiện nay (NASNet, DARTS). AutoShot theo triết lý cell-based: không gian tìm kiếm định nghĩa ở cấp block (DDCNNV2/A/B/C), toàn mạng là tổ hợp 6 block độc lập. Bảng~\ref{tab:search_space_types} tổng hợp đặc điểm của ba loại.

\begin{table}[H]
\centering
\caption{Đặc điểm của các loại không gian tìm kiếm NAS}
\label{tab:search_space_types}
\begin{tabularx}{\textwidth}{lYYY}
\toprule
\textbf{Loại không gian} & \textbf{Kích thước} & \textbf{Linh hoạt} & \textbf{Chi phí tìm kiếm} \\
\midrule
Vĩ mô (Macro) & Rất lớn & Rất cao & Rất cao \\
Phân cấp (Hierarchical) & Lớn & Cao & Cao \\
Dựa trên tế bào (Cell-based) & Trung bình & Trung bình & Trung bình \\
\midrule
AutoShot & $83.9$M kiến trúc & Cao (4 loại block) & Trung bình (SuperNet) \\
\bottomrule
\end{tabularx}
\end{table}


\subsection{Chiến lược tìm kiếm}
\label{subsec:search_strategy}

Chiến lược tìm kiếm khám phá không gian $\mathcal{A}$ cân bằng giữa \textit{khám phá} vùng chưa biết và \textit{khai thác} vùng đã biết tốt \cite{white2023_nas1000}.

\subsubsection{Kỹ thuật tối ưu hóa hộp đen}
\label{subsubsec:blackbox}

Các kỹ thuật hộp đen huấn luyện từng kiến trúc độc lập và không yêu cầu gradient, phù hợp với không gian rời rạc. Ba hướng chính: \textit{RL} \cite{elsken2018_nas_survey} mô hình hóa tìm kiếm như bài toán quyết định tuần tự với RNN làm controller, nhận F1 validation làm phần thưởng — chi phí cao (hàng trăm đến hàng nghìn GPU-giờ); \textit{EA} \cite{elsken2018_nas_survey} tiến hóa quần thể kiến trúc qua đột biến và lai ghép — dễ song song hóa nhưng vẫn tốn đánh giá; \textit{Tối ưu hóa Bayesian (BO)} \cite{elsken2018_nas_survey} xây dựng Gaussian Process xấp xỉ hàm hiệu suất và dùng hàm thu thập để chọn ứng viên tiếp theo một cách thông minh, hiệu quả mẫu nhất trong nhóm. AutoShot sử dụng BO với GP và hàm PoF ở giai đoạn tìm kiếm (Mục~\ref{subsec:autoshot_nas}).

\subsubsection{Kỹ thuật một lần bắn và DARTS}
\label{subsubsec:one_shot}

Các phương pháp một lần bắn (One-Shot) mã hóa toàn bộ không gian tìm kiếm vào một \textit{supernetwork} duy nhất, cho phép tất cả kiến trúc con chia sẻ trọng số và được đánh giá nhanh mà không cần huấn luyện riêng. \textbf{DARTS} \cite{ren2020_nas_survey} là biến thể phổ biến nhất, chuyển không gian rời rạc thành dạng liên tục qua softmax có trọng số trên các phép toán ứng viên; tuy giảm chi phí xuống vài GPU-giờ nhưng gặp vấn đề \textit{khoảng cách rời rạc hóa} khi argmax cuối cùng phá vỡ biểu diễn liên tục đã học --- điểm yếu chí mạng với không gian rời rạc như các biến thể DDCNN.

AutoShot vì vậy chọn \textit{Single-Path One-Shot} (SPOS): tại mỗi bước huấn luyện, một kiến trúc con được lấy mẫu đồng đều và chỉ trọng số của đường đi đó được cập nhật. Cách làm này tránh tương quan giữa các nhánh và hoạt động ổn định với không gian hoàn toàn rời rạc; sau khi SuperNet hội tụ, mọi kiến trúc con có thể được đánh giá tức thời bằng cách kế thừa trọng số chia sẻ.

Bảng~\ref{tab:search_strategy_comparison} tóm tắt đặc điểm và đánh đổi của các chiến lược tìm kiếm chính, giúp định vị AutoShot trong toàn cảnh NAS.

\begin{table}[H]
\centering
\caption{So sánh các chiến lược tìm kiếm kiến trúc}
\label{tab:search_strategy_comparison}
\small
\begin{tabularx}{\textwidth}{lYYYY}
\toprule
\textbf{Chiến lược} & \textbf{Chi phí GPU} & \textbf{Yêu cầu gradient} & \textbf{Không gian rời rạc} & \textbf{Ví dụ} \\
\midrule
Học tăng cường (RL) & Rất cao & Không & Có & NASNet \\
Tiến hóa (EA) & Cao & Không & Có & AmoebaNet \\
Tối ưu hóa Bayesian (BO) & Trung bình & Không & Có & AutoShot \\
DARTS (One-Shot) & Thấp & Có & Không & PC-DARTS \\
Single-Path One-Shot & Thấp & Không & Có & SPOS, AutoShot \\
\bottomrule
\end{tabularx}
\end{table}

\subsection{Chiến lược ước tính hiệu suất}
\label{subsec:performance_estimation}

Ước tính hiệu suất thay thế huấn luyện đầy đủ (vài giờ GPU/kiến trúc) bằng xấp xỉ nhanh: (i) \textit{Lower-fidelity} rút ngắn huấn luyện hoặc dùng tập con dữ liệu; (ii) \textit{Zero-cost proxy} chỉ cần một lần lan truyền để gán điểm cho kiến trúc; (iii) \textit{Weight sharing} (chiến lược AutoShot) — SuperNet chia sẻ trọng số cho tất cả kiến trúc con, cho phép đánh giá bất kỳ kiến trúc trong vài giây. Bảng~\ref{tab:perf_estimation} tóm tắt đánh đổi chi phí–độ chính xác.

\begin{table}[H]
\centering
\caption{So sánh các chiến lược ước tính hiệu suất trong NAS}
\label{tab:perf_estimation}
\begin{tabularx}{\textwidth}{lYYY}
\toprule
\textbf{Chiến lược} & \textbf{Chi phí đánh giá} & \textbf{Độ chính xác xếp hạng} & \textbf{Ví dụ ứng dụng} \\
\midrule
Huấn luyện đầy đủ & Rất cao & Cao nhất & NASNet (RL) \\
Lower-fidelity & Trung bình & Trung bình & Hyperband-NAS \\
Zero-cost proxy & Gần bằng 0 & Thấp–Trung bình & NASWOT, GradNorm \\
Weight sharing (one-shot) & Thấp & Trung bình–Cao & SPOS, AutoShot \\
\bottomrule
\end{tabularx}
\end{table}

Trong AutoShot: không gian tìm kiếm = 83.9M kiến trúc (phương trình~\eqref{eq:search_space_size}), chiến lược = SPOS + BO (Mục~\ref{subsec:autoshot_nas}), ước tính hiệu suất = weight sharing từ SuperNet.

Điểm khởi đầu của AutoShot là TransNetV2 — một kiến trúc được thiết kế thủ công chuyên biệt cho SBD. AutoShot giữ nguyên các thành phần cốt lõi (DDCNN, similarity branch, two-head, sliding window) nhưng đa dạng hóa cách kết hợp tích chập không gian–thời gian thành bốn biến thể, tạo không gian tìm kiếm phong phú hơn. Sau đó, NAS tự động tìm cấu hình tối ưu thay vì chuyên gia quyết định thủ công. Phần tiếp theo sẽ phân tích nền tảng TransNetV2 trước khi trình bày những đóng góp đặc trưng của AutoShot.


\section{Nền tảng: Mô hình TransNetV2}
\label{sec:transnetv2}

TransNetV2 là mô hình mạng nơ-ron sâu được Souček và Lokoč đề xuất năm 2020 \cite{soucek2020transnetv2}, với mục tiêu phát hiện ranh giới cảnh quay trong video một cách nhanh chóng và chính xác. Mô hình kế thừa TransNet (2019) và mang lại nhiều cải tiến căn bản về kiến trúc và quy trình huấn luyện. Mã nguồn và trọng số huấn luyện sẵn được tác giả công bố rộng rãi, tạo nền tảng quan trọng cho cộng đồng nghiên cứu và là điểm khởi đầu trực tiếp của AutoShot.

\subsection{Tổng quan và vị trí trong dòng chảy nghiên cứu SBD}
\label{subsec:transnetv2_overview}

Trước TransNetV2, phần lớn các phương pháp SBD dùng đặc trưng thủ công (histogram màu, pixel difference) hoặc mạng CNN xử lý từng cặp khung hình riêng lẻ. Nhóm phương pháp đầu dễ cài đặt nhưng thiếu khả năng nắm bắt ngữ cảnh thị giác phức tạp; nhóm sau học được đặc trưng phong phú hơn nhưng không mô hình hóa ngữ cảnh thời gian dài hạn — quan trọng để phân biệt chuyển cảnh đột ngột (thay đổi đột ngột) với chuyển cảnh dần (thay đổi dần dần như dissolve, fade).

TransNetV2 giải quyết cả hai hạn chế này đồng thời: khối tích chập giãn nở 3D (DDCNN) học đặc trưng thời gian đa tỉ lệ trực tiếp từ dữ liệu thô, trong khi nhánh đặc trưng tương đồng cung cấp tín hiệu bổ sung về sự thay đổi nội dung giữa các khung hình. Tích chập phân tách không gian–thời gian (tích chập phân tách) giảm số tham số xuống còn khoảng 6.5 triệu trong khi vẫn duy trì trường nhìn thời gian rộng lên tới 97 khung hình.

So với TransNetV1, TransNetV2 cải thiện F1-score đáng kể trên các bộ dữ liệu chuẩn: từ 73.5\% lên 77.9\% trên ClipShots và từ 92.9\% lên 96.2\% trên BBC Planet Earth. Những cải thiện này đến từ sự kết hợp đồng bộ của bốn thay đổi quan trọng, được tổng hợp trong Bảng~\ref{tab:transnetv1_v2}.

\begin{table}[H]
\centering
\caption{Các cải tiến chính từ TransNetV1 sang TransNetV2}
\label{tab:transnetv1_v2}
\begin{tabularx}{\textwidth}{L{3.3cm}QQ}
\toprule
\textbf{Thành phần} & \textbf{TransNetV1} & \textbf{TransNetV2} \\
\midrule
Chuẩn hóa & Không có & BatchNorm sau mỗi tích chập \\
Kết nối & Tuần tự & Skip connection sau mỗi 2 khối \\
Kernel tích chập & 3D đầy đủ & Phân tách không gian–thời gian \\
Nhánh bổ sung & Không có & Learned + Histogram similarity \\
Đầu ra & Một head & Hai heads (single + all frame) \\
Dữ liệu huấn luyện & Nhân tạo 100\% & Nhân tạo 85\% + Thực 15\% \\
\midrule
F1 ClipShots & 73.5\% & 77.9\% \\
F1 BBC Planet Earth & 92.9\% & 96.2\% \\
\bottomrule
\end{tabularx}
\end{table}

Mỗi thay đổi trong Bảng~\ref{tab:transnetv1_v2} giải quyết một vấn đề cụ thể. BatchNorm và skip connection ổn định quá trình huấn luyện khi mô hình sâu hơn. Phân tách kernel giảm nguy cơ quá khớp trên dữ liệu nhân tạo có phân phối hạn chế. Nhánh similarity cung cấp tín hiệu bổ sung không thể học được chỉ từ tích chập cục bộ. Hai heads buộc mô hình phân biệt giữa "tại đâu là ranh giới" và "vùng nào là chuyển cảnh", hai câu hỏi bổ sung cho nhau. Cuối cùng, bổ sung dữ liệu thực vào 15\% tập huấn luyện thu hẹp khoảng cách phân phối giữa huấn luyện và kiểm thử.

\subsection{Kiến trúc mô hình}
\label{subsec:transnetv2_arch}

TransNetV2 xử lý chuỗi khung hình đầu vào có kích thước $48 \times 27 \times 3$ pixel với độ dài 100 khung, trả ra xác suất ranh giới cho mỗi khung hình. Kiến trúc gồm hai nhánh song song — nhánh DDCNN chính và nhánh đặc trưng tương đồng — được kết hợp trước các đầu phân loại, như minh họa trong Hình~\ref{fig:transnetv2_arch}.

\begin{figure}[H]
    \centering
    \resizebox{0.9\textwidth}{!}{
    \begin{tikzpicture}[>=Latex,font=\footnotesize]
        \node[draw, rounded corners, fill=blue!10, text width=2.7cm, minimum height=1.0cm, align=center] (input) at (0,0) {Input clip\\$100 \times 48 \times 27 \times 3$};
        \node[draw, rounded corners, fill=green!12, text width=3.5cm, minimum height=1.1cm, align=center] (main) at (4.4,1.1) {Nhánh chính\\6 khối DDCNN giãn nở\\(+ skip/pool mỗi 2 khối)};
        \node[draw, rounded corners, fill=yellow!20, text width=2.8cm, minimum height=1.0cm, align=center] (hist) at (4.4,-1.1) {Histogram\\similarity RGB};
        \node[draw, rounded corners, fill=yellow!20, text width=2.8cm, minimum height=1.0cm, align=center] (learn) at (8.0,-1.1) {Learned\\similarity};
        \node[draw, rounded corners, fill=orange!18, text width=2.4cm, minimum height=1.0cm, align=center] (merge) at (8.0,1.1) {Fusion\\đặc trưng};
        \node[draw, rounded corners, fill=purple!12, text width=2.4cm, minimum height=1.0cm, align=center] (head1) at (11.7,2.1) {Single-frame\\head};
        \node[draw, rounded corners, fill=purple!12, text width=2.4cm, minimum height=1.0cm, align=center] (head2) at (11.7,0.1) {All-frame\\head};

        \draw[->, thick] (input) -- (main);
        \draw[->, thick] (input) -- (hist);
        \draw[->, thick] (hist) -- (learn);
        \draw[->, thick] (learn) -- (merge);
        \draw[->, thick] (main) -- (merge);
        \draw[->, thick] (merge) -- (head1);
        \draw[->, thick] (merge) -- (head2);
        \node[align=left] at (14.35,1.1) {\footnotesize $p_t^{(1)}$, $p_t^{(\mathrm{all})}$\\\footnotesize xác suất ranh giới theo khung};
    \end{tikzpicture}
    }
    \caption[Tổng quan kiến trúc TransNetV2]{Tổng quan kiến trúc TransNetV2: 6 khối DDCNN giãn nở xếp chồng (nhánh chính) kết hợp với nhánh đặc trưng tương đồng (histogram RGB + learned similarity), đầu ra qua hai đầu phân loại song song (single-frame head và all-frame head). Kết nối tắt và pooling trung bình không gian được áp dụng sau mỗi hai khối DDCNN.}
    \label{fig:transnetv2_arch}
\end{figure}

\subsubsection{Khối tích chập giãn nở DDCNN}
\label{subsubsec:ddcnn}

Xương sống của TransNetV2 là 6 khối tích chập giãn nở (Dilated DCNN — DDCNN) xếp chồng tuần tự. Mỗi khối DDCNN gồm 4 nhánh tích chập song song với độ giãn nở (dilation) theo chiều thời gian tăng dần theo cấp số nhân: $d \in \{1, 2, 4, 8\}$. Đầu ra của 4 nhánh được nối ghép (concatenate) theo chiều kênh, tiếp theo là BatchNorm và ReLU. Sau mỗi hai khối DDCNN, một skip connection và lớp pooling trung bình không gian được thêm vào để ổn định luồng gradient và giảm dần độ phân giải không gian.

Hình~\ref{fig:ddcnn_block} minh họa cấu trúc chi tiết của một khối DDCNN với 4 nhánh dilation song song.

\begin{figure}[H]
    \centering
    \resizebox{0.85\textwidth}{!}{
    \begin{tikzpicture}[>=Latex,font=\scriptsize]
        \node[draw, rounded corners, fill=blue!10, text width=2.3cm, minimum height=0.95cm, align=center] (input) at (0,0) {Input feature\\$X_t$};
        \node[draw, rounded corners, fill=gray!10, text width=1.7cm, minimum height=0.9cm, align=center] (s1) at (3.1,2.1) {Spatial\\$1\times k\times k$};
        \node[draw, rounded corners, fill=gray!10, text width=1.7cm, minimum height=0.9cm, align=center] (s2) at (3.1,0.7) {Spatial\\$1\times k\times k$};
        \node[draw, rounded corners, fill=gray!10, text width=1.7cm, minimum height=0.9cm, align=center] (s3) at (3.1,-0.7) {Spatial\\$1\times k\times k$};
        \node[draw, rounded corners, fill=gray!10, text width=1.7cm, minimum height=0.9cm, align=center] (s4) at (3.1,-2.1) {Spatial\\$1\times k\times k$};

        \node[draw, rounded corners, fill=green!12, text width=2.2cm, minimum height=1.0cm, align=center] (t1) at (6.3,2.1) {Temporal\\$k\times1\times1$\\$d=1$};
        \node[draw, rounded corners, fill=green!12, text width=2.2cm, minimum height=1.0cm, align=center] (t2) at (6.3,0.7) {Temporal\\$k\times1\times1$\\$d=2$};
        \node[draw, rounded corners, fill=green!12, text width=2.2cm, minimum height=1.0cm, align=center] (t3) at (6.3,-0.7) {Temporal\\$k\times1\times1$\\$d=4$};
        \node[draw, rounded corners, fill=green!12, text width=2.2cm, minimum height=1.0cm, align=center] (t4) at (6.3,-2.1) {Temporal\\$k\times1\times1$\\$d=8$};

        \node[draw, rounded corners, fill=orange!18, text width=1.9cm, minimum height=1.1cm, align=center] (concat) at (9.3,0) {Concat\\theo kênh};
        \node[draw, rounded corners, fill=purple!12, text width=2.0cm, minimum height=1.1cm, align=center] (bn) at (12.0,0) {BatchNorm\\+ ReLU};
        \node[draw, rounded corners, fill=blue!10, text width=1.9cm, minimum height=1.1cm, align=center] (out) at (14.6,0) {Output\\$Y_t$};

        \draw[->, thick] (input) -- (s1);
        \draw[->, thick] (input) -- (s2);
        \draw[->, thick] (input) -- (s3);
        \draw[->, thick] (input) -- (s4);
        \draw[->, thick] (s1) -- (t1);
        \draw[->, thick] (s2) -- (t2);
        \draw[->, thick] (s3) -- (t3);
        \draw[->, thick] (s4) -- (t4);
        \draw[->, thick] (t1) -- (concat);
        \draw[->, thick] (t2) -- (concat);
        \draw[->, thick] (t3) -- (concat);
        \draw[->, thick] (t4) -- (concat);
        \draw[->, thick] (concat) -- (bn);
        \draw[->, thick] (bn) -- (out);
    \end{tikzpicture}
    }
    \caption[Cấu trúc khối DDCNN trong TransNetV2]{Cấu trúc chi tiết của khối DDCNN trong TransNetV2: 4 nhánh tích chập thời gian 1D với độ giãn nở $d \in \{1,2,4,8\}$ song song nhau, đầu ra được nối ghép theo chiều kênh. Mỗi nhánh trước tiên qua tích chập không gian 2D ($1\times k\times k$) rồi mới qua tích chập thời gian 1D giãn nở ($k\times 1\times 1$).}
    \label{fig:ddcnn_block}
\end{figure}

Nhờ cấu trúc giãn nở, trường nhìn thời gian (temporal receptive field) tại khối DDCNN thứ $k$ có thể được tính theo quy nạp. Với kernel thời gian kích thước $K = 3$ và 4 nhánh dilation $\{1, 2, 4, 8\}$, trường nhìn cộng dồn sau $k$ khối là:
\begin{equation}
    R_k = R_{k-1} + 2 \cdot (K - 1) \cdot \sum_{d \in \{1,2,4,8\}} d
    = R_{k-1} + 2 \cdot 2 \cdot 15 = R_{k-1} + 60,
    \label{eq:receptive_field}
\end{equation}
với $R_0 = 1$ (một khung hình đơn lẻ). Cơ chế cộng dồn cho $R_1 = 61$, $R_2 = 121$, \dots, $R_6 = 361$ trong trường hợp lý tưởng. Tuy nhiên, chiều dài chuỗi đầu vào trong huấn luyện và suy diễn được cố định ở mức 100 khung hình, nên trường nhìn bị cắt cụt; kết hợp với các lớp pooling xen kẽ, trường nhìn hiệu quả (effective receptive field) tập trung vào khoảng 97 khung hình — đủ bao quát toàn bộ chuỗi đầu vào, giúp mô hình nắm bắt ngữ cảnh dài hạn để phân biệt cảnh tan dần (dissolve) với chuyển động chậm (slow motion).

\subsubsection{Phân tách tích chập không gian–thời gian}
\label{subsubsec:factorized_conv}

Thay vì dùng kernel tích chập 3D đầy đủ kích thước $k \times k \times k$, TransNetV2 phân tách mỗi phép tích chập thành hai bước tuần tự nhằm giảm số tham số và buộc mạng học riêng biệt hai loại thông tin:
\begin{enumerate}
    \item \textbf{Tích chập không gian 2D} với kernel $1 \times k \times k$: xử lý từng khung hình trong chuỗi một cách độc lập theo hai chiều không gian (height, width), học đặc trưng nội dung thị giác như cạnh, kết cấu, vùng màu.
    \item \textbf{Tích chập thời gian 1D giãn nở} với kernel $k \times 1 \times 1$ và độ giãn nở $d$: xử lý chuỗi đặc trưng theo trục thời gian, học sự biến đổi chuyển động và chuyển cảnh.
\end{enumerate}

Phân tách giảm số tham số còn $\approx 44\%$ so với kernel 3D đầy đủ: với $k=3$ và $C_{in}=C_{out}=C_{mid}=C$, tỉ lệ là $(9C^2+3C^2)/(27C^2)=12/27\approx44\%$, tức tiết kiệm $\approx56\%$ tham số bất kể số kênh. Ngoài lợi ích về bộ nhớ, việc ép xử lý không gian trước rồi mới xử lý thời gian tạo ra thiên kiến quy nạp phù hợp với SBD: nội dung khung thay đổi theo không gian, còn chuyển cảnh là hiện tượng theo thời gian — hai biến đổi nên được mô hình hóa riêng biệt.

\subsubsection{Nhánh đặc trưng tương đồng và hai đầu phân loại}
\label{subsubsec:similarity_branch}

Song song với DDCNN, TransNetV2 tính ma trận tương đồng qua hai nguồn bổ sung: \textit{learned similarity} (cosine similarity giữa đặc trưng $\ell_2$-chuẩn hóa trích từ $L_2,L_4,L_6$) và \textit{histogram RGB similarity} (cosine distance của histogram 512-bin). Hai loại tương đồng (mỗi loại 101 chiều) được ghép thành 202-D, qua Dense(128)+ReLU rồi kết hợp với đầu ra DDCNN. Learned similarity nhạy với thay đổi ngữ nghĩa; histogram nhạy với thay đổi màu sắc thô — sự kết hợp tăng tính bền vững qua nhiều loại chuyển cảnh.

\textbf{Hai đầu phân loại:} \textit{Single-frame head} dự đoán xác suất tại đúng một khung trung tâm của ranh giới (đầu ra chính khi suy diễn); \textit{All-frame head} dự đoán xác suất "thuộc về vùng chuyển cảnh" (hỗ trợ học phạm vi chuyển cảnh dần). Cơ chế hai đầu là cải tiến then chốt so với TransNetV1.

\subsection{Hàm mục tiêu và quy trình huấn luyện}
\label{subsec:transnetv2_training}

\subsubsection{Hàm mất mát weighted cross-entropy}

TransNetV2 tối thiểu hóa weighted cross-entropy tổng hợp hai nhánh kèm $\ell_2$:
\begin{equation}
    \mathcal{L}(\theta)
    = \frac{\lambda_1}{KM} \sum_{i,t} \mathrm{CE}\!\bigl(\hat{y}^{(1)}_{i,t},\, y^{(1)}_{i,t}\bigr)
    + \frac{\lambda_2}{KM} \sum_{i,t} \mathrm{CE}\!\bigl(\hat{y}^{(2)}_{i,t},\, y^{(2)}_{i,t}\bigr)
    + \lambda_{\mathrm{reg}} \|\theta\|^2,
    \label{eq:transnetv2_loss}
\end{equation}
với $\lambda_1=5.0$, $\lambda_2=0.1$, $\lambda_{\mathrm{reg}}=10^{-4}$. Trọng số $\lambda_1 \gg \lambda_2$ ưu tiên nhánh single-frame vốn mất cân bằng lớp nghiêm trọng hơn.

\subsubsection{Chiến lược dữ liệu và tăng cường}
\label{subsubsec:data_strategy}

Tập huấn luyện kết hợp 85\% dữ liệu nhân tạo (chuyển cảnh đột ngột và dissolve tạo từ IACC.3 + ClipShots) và 15\% dữ liệu thực (ClipShots). Tăng cường dữ liệu được áp dụng \textit{nhất quán cho toàn chuỗi} (flip, màu sắc, color transfer 10\%) để tránh artifact thời gian giả tạo. Chuyển đổi màu giả lập trường hợp khó: hai cảnh liền kề có tông màu tương đồng, buộc mô hình học đặc trưng cấu trúc thay vì chỉ dựa vào màu sắc (Bảng~\ref{tab:transnetv2_data}).

\begin{table}[H]
\centering
\caption{Chiến lược dữ liệu huấn luyện của TransNetV2}
\label{tab:transnetv2_data}
\begin{tabularx}{\textwidth}{lYr}
\toprule
\textbf{Loại mẫu} & \textbf{Nguồn / Cách tạo} & \textbf{Tỉ lệ} \\
\midrule
Chuỗi thực & ClipShots (ranh giới cảnh thực) & 15\% \\
chuyển cảnh đột ngột nhân tạo & IACC.3 + ClipShots, nối 2 cảnh & 35\% \\
Dissolve nhân tạo & IACC.3 + ClipShots, hiệu ứng hòa tan & 50\% \\
\midrule
Tăng cường & Flip, màu sắc, color transfer (10\%) & 100\% \\
\bottomrule
\end{tabularx}
\end{table}

\subsection{Thuật toán suy diễn với cửa sổ trượt}
\label{subsec:transnetv2_algorithms}

Thuật toán~\ref{alg:transnetv2_inference} mô tả cơ chế cửa sổ trượt xử lý video có độ dài tùy ý. TransNetV2 chỉ nhận đầu vào có độ dài cố định 100 khung, nên cần chia video dài thành các đoạn chồng lấp. Mỗi cửa sổ 100 khung di chuyển với bước nhảy 50 khung; chỉ 50 khung hình ở giữa (vị trí 25–74) được giữ lại trong kết quả, tránh hiệu ứng biên ở hai đầu cửa sổ nơi ngữ cảnh thời gian bị thiếu.

\begin{algorithm}[H]
\caption{TransNetV2 Inference}
\label{alg:transnetv2_inference}
\begin{algorithmic}[1]
\Function{Predict}{$video,\; \theta^*,\; \tau = 0.5$}
    \State $frames \gets \text{ExtractFrames}(video)$
    \State $frames \gets \text{Resize}(frames,\; 48 \times 27 \times 3)$
    \State $predictions \gets []$
    \For{$start = 0 \text{ to } |frames| - 100 \text{ step } 50$}
        \State $\hat{y}^{(1)} \gets \text{Forward}(\theta^*,\; frames[start : start+100])$
        \State $predictions.\text{extend}(\hat{y}^{(1)}[25:75])$
        \Comment{Keep middle 50 frames}
    \EndFor
    \State $\text{positions} \gets \{t : predictions[t] > \tau\}$
    \State $boundaries \gets \text{GroupConsecutive}(\text{positions})$
    \Comment{Gradual: multiple consecutive frames}
    \State \Return $boundaries$
\EndFunction
\end{algorithmic}
\end{algorithm}

Quy trình suy diễn này được AutoShot tái sử dụng nguyên vẹn: chỉ 50 khung giữa được giữ lại, đảm bảo video được bao phủ đầy đủ và tránh sai số ở biên cửa sổ.

\subsection{Thông số cấu hình}
\label{subsec:transnetv2_complexity}

Bảng~\ref{tab:transnetv2_params} tổng hợp các siêu tham số chính, phục vụ đối chiếu với AutoShot ở Mục~\ref{sec:comparison}. Mô hình $\approx6.5$M tham số, huấn luyện $\approx17$ giờ trên V100.

\begin{table}[H]
\centering
\caption{Các thông số cấu hình của TransNetV2}
\label{tab:transnetv2_params}
\begin{tabularx}{\textwidth}{lY}
\toprule
\textbf{Tham số} & \textbf{Giá trị} \\
\midrule
Độ phân giải đầu vào & $48 \times 27 \times 3$ \\
Độ dài chuỗi huấn luyện & 100 khung hình \\
Bước trượt suy diễn & 50 khung hình \\
Kích thước batch & 16 \\
Tốc độ học & 0.01 \\
Momentum (SGD) & 0.9 \\
Số epoch huấn luyện & 50 \\
Trọng số loss $(\lambda_1, \lambda_2)$ & 5.0, 0.1 \\
Chính quy hóa $(\lambda_{\mathrm{reg}})$ & $10^{-4}$ \\
Trường nhìn thời gian hiệu quả & $\approx$ 97 khung hình \\
Tổng số tham số & $\approx$ 6.5 triệu \\
Bộ nhớ GPU (batch 16) & $\approx$ 2 GB \\
Thời gian huấn luyện (V100) & $\approx$ 17 giờ \\
\bottomrule
\end{tabularx}
\end{table}

\noindent
Từ nền tảng NAS và TransNetV2 đã trình bày, các mục tiếp theo cụ thể hóa khung NAS thành hai mô hình. Trước hết, Mục~\ref{sec:autoshot} mô tả AutoShot --- một instance của khung NAS với 83.9 triệu kiến trúc ứng viên, dùng SuperNet kết hợp Bayesian Optimization để chọn mạng xương sống phù hợp với phân phối video ngắn của SHOT. Tiếp theo, Mục~\ref{sec:autoshotv2} trình bày đóng góp của khóa luận: biến thể AutoShotV2 đóng băng mạng xương sống đã được NAS tối ưu, chỉ huấn luyện lại tầng phân loại với Focal Loss, nhãn many-hot, Gaussian smoothing và temperature scaling --- nhằm cô lập đóng góp của tầng đầu ra so với mạng xương sống.

\section{Mô hình AutoShot — nền tảng kế thừa từ TransNetV2}
\label{sec:autoshot}

\subsection{Động lực kế thừa từ TransNetV2}
\label{subsec:autoshot_motivation}

AutoShot~\cite{zhu2023autoshot} bắt đầu từ một nhận xét thực nghiệm: TransNetV2~\cite{soucek2020transnetv2} là mô hình cơ sở mạnh cho SBD, nhưng chỉ đạt F1 = 79.9\% trên SHOT, thấp hơn kỳ vọng khi chuyển sang miền video ngắn. Vấn đề không nằm ở toàn bộ quy trình TransNetV2, vì các quyết định như cửa sổ trượt, nhánh similarity và hai đầu phân loại vẫn phù hợp. Điểm cần tối ưu là \textbf{backbone}: TransNetV2 dùng cùng một dạng DDCNN cho 6 vị trí, trong khi video ngắn có nhịp cắt nhanh, hiệu ứng chỉnh màu dày đặc và phân phối chuyển cảnh khác video dài.

Vì vậy, AutoShot không thay thế TransNetV2 từ đầu. Mô hình giữ lại quy trình đã được kiểm chứng, rồi dùng Neural Architecture Search (NAS) để chọn cấu hình mạng xương sống phù hợp hơn với SHOT và ClipShots. Cách tổ chức này giúp so sánh công bằng: nếu hiệu năng tăng, nguyên nhân chính đến từ cơ chế chọn kiến trúc bằng NAS chứ không phải từ việc đổi toàn bộ quy trình suy diễn hay hàm mục tiêu.

Hai luận điểm kỹ thuật được kiểm chứng trong AutoShot:
\begin{itemize}
    \item \textbf{H1}: mạng xương sống cố định với 6 khối DDCNN đồng nhất chưa tối ưu cho video ngắn; chuyên biệt hóa theo từng vị trí có thể cải thiện F1.
    \item \textbf{H2}: NAS one-shot trên không gian gồm các biến thể DDCNN và tùy chọn Transformer có thể tìm kiến trúc vượt TransNetV2 với chi phí huấn luyện chấp nhận được.
\end{itemize}

Bảng~\ref{tab:autoshot_inherited} phân rõ thành phần nào của AutoShot được kế thừa từ TransNetV2 và thành phần nào là đóng góp của công trình AutoShot gốc.

\begin{table}[H]
\centering
\caption{Thành phần kế thừa từ TransNetV2 và đóng góp của công trình AutoShot gốc}
\label{tab:autoshot_inherited}
\begin{tabularx}{\textwidth}{llY}
\toprule
\textbf{Thành phần} & \textbf{Kế thừa / Mới} & \textbf{Mô tả} \\
\midrule
Độ phân giải đầu vào ($48\times27$) & Kế thừa & Giữ nguyên \\
Nhánh similarity (learned + hist.) & Kế thừa & Giữ nguyên \\
Hai đầu phân loại (single + all) & Kế thừa & Giữ nguyên \\
Hàm mất mát (WCE hai nhánh) & Kế thừa & Giữ nguyên \\
Suy diễn sliding window & Kế thừa & Giữ nguyên \\
\midrule
Loại khối mạng xương sống & \textbf{Mới} & 4 biến thể DDCNNV2/A/B/C \\
Khối Transformer (vị trí 7) & \textbf{Mới} & Tùy chọn self-attention \\
Quy trình NAS (SuperNet + BO) & \textbf{Mới} & Tìm kiếm kiến trúc tự động \\
chưng cất tri thức & \textbf{Mới} & Học từ nhãn mềm \\
Weight Grafting & \textbf{Mới} & Kết hợp trọng số theo entropy \\
\bottomrule
\end{tabularx}
\end{table}

\subsection{Không gian tìm kiếm kiến trúc}
\label{subsec:autoshot_search_space}

Không gian tìm kiếm của AutoShot được thiết kế hẹp có chủ đích: chỉ thay đổi các khối trích xuất đặc trưng trong mạng xương sống, còn phần vào/ra và quy trình suy diễn vẫn giữ theo TransNetV2. Cụ thể, mô hình có 7 quyết định kiến trúc nối tiếp: 6 vị trí đầu chọn khối tích chập phân tách 3D, vị trí cuối chọn số lớp Transformer 1D. Hình~\ref{fig:autoshot_overall} minh họa sơ đồ tổng thể, trong đó mỗi biến $\alpha_i$ là một quyết định rời rạc của NAS.

\begin{figure}[H]
    \centering
    \resizebox{\textwidth}{!}{
    \begin{tikzpicture}[>=Latex,font=\footnotesize, node distance=0.35cm]
        \tikzset{
            iopad/.style={draw, rounded corners, fill=blue!10, text width=2.0cm, minimum height=1.15cm, align=center},
            nasblock/.style={draw, thick, rounded corners, fill=green!15, text width=2.0cm, minimum height=1.15cm, align=center},
            nastrans/.style={draw, thick, rounded corners, fill=orange!20, text width=2.0cm, minimum height=1.15cm, align=center},
            shared/.style={draw, rounded corners, fill=purple!12, text width=2.0cm, minimum height=1.15cm, align=center}
        }
        \node[iopad] (in) at (0,0) {Video\\$48{\times}27{\times}3$\\(100 frames)};
        \node[nasblock, right=of in] (b1) {Block 1\\$\alpha_1 \in \mathcal{V}$};
        \node[nasblock, right=of b1] (b2) {Block 2\\$\alpha_2 \in \mathcal{V}$};
        \node[nasblock, right=of b2] (b3) {Block 3\\$\alpha_3 \in \mathcal{V}$};
        \node[nasblock, right=of b3] (b4) {Block 4\\$\alpha_4 \in \mathcal{V}$};
        \node[nasblock, right=of b4] (b5) {Block 5\\$\alpha_5 \in \mathcal{V}$};
        \node[nasblock, right=of b5] (b6) {Block 6\\$\alpha_6 \in \mathcal{V}$};
        \node[nastrans, right=of b6] (t7) {Transformer\\$\alpha_7 \in \{0..4\}$};
        \node[shared, right=of t7] (head) {Similarity branch\\+ 2 heads};
        \node[iopad, right=of head] (out) {$\hat{y}^{(1)}, \hat{y}^{(2)}$};

        \draw[->, thick] (in) -- (b1);
        \draw[->, thick] (b1) -- (b2);
        \draw[->, thick] (b2) -- (b3);
        \draw[->, thick] (b3) -- (b4);
        \draw[->, thick] (b4) -- (b5);
        \draw[->, thick] (b5) -- (b6);
        \draw[->, thick] (b6) -- (t7);
        \draw[->, thick] (t7) -- (head);
        \draw[->, thick] (head) -- (out);

        \draw[decorate, decoration={brace, amplitude=7pt, mirror}, thick]
            ([yshift=-10pt]b1.south west) -- ([yshift=-10pt]b6.south east)
            node[midway, below=9pt, align=center] {\footnotesize 6 khối tích chập (NAS search):\\ $\mathcal{V} = \{\mathrm{DDCNNV2}, \mathrm{V2A}, \mathrm{V2B}, \mathrm{V2C}\}$ (16 cấu hình/khối)};
        \draw[decorate, decoration={brace, amplitude=7pt, mirror}, thick]
            ([yshift=-10pt]t7.south west) -- ([yshift=-10pt]t7.south east)
            node[midway, below=9pt, align=center] {\footnotesize Tự chú ý\\(0--4 lớp)};
        \draw[decorate, decoration={brace, amplitude=7pt, mirror}, thick]
            ([yshift=-10pt]head.south west) -- ([yshift=-10pt]head.south east)
            node[midway, below=9pt, align=center] {\footnotesize Kế thừa\\TransNetV2};
    \end{tikzpicture}
    }
    \caption[Sơ đồ tổng thể mô hình AutoShot]{Sơ đồ tổng thể mô hình AutoShot gồm 7 vị trí kiến trúc cần tìm kiếm. Sáu vị trí đầu là các khối tích chập phân tách 3D chọn từ $\mathcal{V} = \{\mathrm{DDCNNV2}, \mathrm{DDCNNV2A}, \mathrm{DDCNNV2B}, \mathrm{DDCNNV2C}\}$ với các siêu tham số $(nc, nd)$; vị trí thứ 7 là khối Transformer 1D với số lớp tự chú ý $\alpha_7 \in \{0, 1, 2, 3, 4\}$. Nhánh similarity và hai đầu phân loại được kế thừa nguyên vẹn từ TransNetV2, do đó mọi cải thiện đo được đều quy về sự khác biệt ở 7 quyết định $\alpha_1, \dots, \alpha_7$ của NAS.}
    \label{fig:autoshot_overall}
\end{figure}

\subsubsection{Bốn biến thể khối tích chập}
\label{subsubsec:four_variants}

Hình~\ref{fig:autoshot_blocks} minh họa bốn biến thể khối tích chập 3D phân tách trong không gian tìm kiếm. Tất cả đều kế thừa nguyên tắc tích chập phân tách từ TransNetV2 (Mục~\ref{subsubsec:factorized_conv}), nhưng khác nhau ở cách chia sẻ spatial convolution và cách ghép đặc trưng spatial--temporal. Bảng~\ref{tab:block_variants} tóm tắt vai trò của từng biến thể.

\begin{figure}[H]
    \centering
    \resizebox{0.98\textwidth}{!}{
    \begin{tikzpicture}[>=Latex,font=\scriptsize]
        \tikzset{
            titlebox/.style={draw, rounded corners, fill=gray!15, minimum width=3.8cm, minimum height=0.75cm, align=center, font=\bfseries\scriptsize},
            data/.style={draw, rounded corners, fill=blue!10, minimum width=3.2cm, minimum height=0.8cm, align=center},
            spatial/.style={draw, rounded corners, fill=green!12, minimum width=3.2cm, minimum height=0.8cm, align=center},
            temporal/.style={draw, rounded corners, fill=orange!18, minimum width=3.2cm, minimum height=0.8cm, align=center},
            merge/.style={draw, rounded corners, fill=purple!12, minimum width=3.2cm, minimum height=0.8cm, align=center}
        }

        % DDCNNV2
        \begin{scope}[xshift=0cm]
            \node[titlebox] (t1) at (0,3.3) {DDCNNV2};
            \node[data] (i1) at (0,2.2) {Input};
            \node[spatial] (s1) at (0,1.1) {4 nhánh spatial riêng biệt};
            \node[temporal] (p1) at (0,0.0) {4 nhánh temporal giãn nở};
            \node[merge] (m1) at (0,-1.1) {Concat};
            \draw[->, thick] (i1) -- (s1);
            \draw[->, thick] (s1) -- (p1);
            \draw[->, thick] (p1) -- (m1);
        \end{scope}

        % DDCNNV2A
        \begin{scope}[xshift=5.0cm]
            \node[titlebox] (t2) at (0,3.3) {DDCNNV2A};
            \node[data] (i2) at (0,2.2) {Input};
            \node[spatial] (s2) at (0,1.1) {Shared spatial conv};
            \node[temporal] (p2) at (0,0.0) {4 nhánh temporal giãn nở};
            \node[merge] (m2) at (0,-1.1) {Concat};
            \draw[->, thick] (i2) -- (s2);
            \draw[->, thick] (s2) -- (p2);
            \draw[->, thick] (p2) -- (m2);
        \end{scope}

        % DDCNNV2B
        \begin{scope}[xshift=10.0cm]
            \node[titlebox] (t3) at (0,3.3) {DDCNNV2B};
            \node[data] (i3) at (0,2.2) {Input};
            \node[spatial, minimum width=1.8cm] (s3) at (-1.1,0.9) {Spatial};
            \node[temporal, minimum width=1.8cm] (p3) at (1.1,0.9) {Temporal};
            \node[merge] (m3) at (0,-0.3) {Add};
            \draw[->, thick] (i3) -- (s3);
            \draw[->, thick] (i3) -- (p3);
            \draw[->, thick] (s3) -- (m3);
            \draw[->, thick] (p3) -- (m3);
        \end{scope}

        % DDCNNV2C
        \begin{scope}[xshift=15.0cm]
            \node[titlebox] (t4) at (0,3.3) {DDCNNV2C};
            \node[data] (i4) at (0,2.2) {Input};
            \node[spatial] (s4) at (0,1.1) {Shared spatial conv};
            \node[temporal] (p4) at (0,0.0) {Temporal giãn nở};
            \node[merge] (m4) at (0,-1.1) {Add};
            \draw[->, thick] (i4) -- (s4);
            \draw[->, thick] (s4) -- (p4);
            \draw[->, thick] (p4) -- (m4);
            \draw[->, thick, bend left=35] (s4.east) to (m4.east);
        \end{scope}
    \end{tikzpicture}
    }
    \caption[Bốn biến thể khối tích chập 3D trong AutoShot]{Bốn biến thể khối tích chập 3D phân tách trong không gian tìm kiếm của AutoShot. Từ trái qua phải: DDCNNV2 (spatial riêng biệt từng nhánh, concat), DDCNNV2A (shared spatial, concat), DDCNNV2B (spatial + temporal song song, cộng dồn theo P3D), DDCNNV2C (spatial truyền tuần tự vào temporal, cộng dồn). Tất cả đều dùng tích chập thời gian 1D giãn nở đa nhánh.}
    \label{fig:autoshot_blocks}
\end{figure}

\begin{table}[H]
\centering
\caption{So sánh bốn biến thể khối tích chập 3D phân tách của AutoShot}
\label{tab:block_variants}
\begin{tabularx}{\textwidth}{lYcY}
\toprule
\textbf{Biến thể} & \textbf{Spatial conv} & \textbf{Phép kết hợp} & \textbf{Nguồn cảm hứng} \\
\midrule
DDCNNV2   & Riêng biệt từng nhánh & Concat & TransNetV2 gốc \\
DDCNNV2A  & Dùng chung (shared)   & Concat & Chia sẻ đặc trưng không gian \\
DDCNNV2B  & Song song với temporal & Add   & Pseudo-3D (P3D) \\
DDCNNV2C  & Tuần tự vào temporal  & Add   & Kết hợp DDCNNV2A + B \\
\bottomrule
\end{tabularx}
\end{table}

Bảng~\ref{tab:block_variants} là nguồn so sánh chính giữa bốn biến thể; điểm cốt lõi cần ghi nhớ là không gian tìm kiếm phủ hai trục thiết kế: \textit{cách chia sẻ spatial convolution} (riêng biệt như V2 vs. dùng chung như V2A, V2C) và \textit{cách ghép spatial--temporal} (Concat như V2/V2A vs. Add lấy cảm hứng P3D như V2B/V2C). Đưa cả hai cơ chế Concat và Add vào không gian tìm kiếm cho phép NAS tự cân nhắc giữa tính phân tách thông tin (Concat) và sự ổn định gradient kiểu ResNet (Add) phù hợp nhất với từng vị trí trong mạng xương sống của video ngắn.

\subsubsection{Kích thước không gian tìm kiếm}
\label{subsubsec:search_space_size}

Với 4 loại khối và các tùy chọn siêu tham số như mô tả trong Bảng~\ref{tab:block_variants}, số cấu hình khả thi cho mỗi khối tích chập là:
\begin{itemize}
    \item DDCNNV2: $3 \times 2 = 6$ cấu hình (3 giá trị $nc$, 2 giá trị $nd$)
    \item DDCNNV2A: $3 \times 2 = 6$ cấu hình
    \item DDCNNV2B: $1 \times 2 = 2$ cấu hình ($nc$ cố định, 2 giá trị $nd$)
    \item DDCNNV2C: $1 \times 2 = 2$ cấu hình
    \item Tổng mỗi khối tích chập: $6 + 6 + 2 + 2 = 16$ cấu hình
\end{itemize}

Khối Transformer thứ 7 có 5 tùy chọn về số lớp self-attention: $\{0, 1, 2, 3, 4\}$. Các 6 khối tích chập được cấu hình độc lập, nên tổng số kiến trúc trong không gian tìm kiếm là:
\begin{equation}
    |\mathcal{A}| = 16^6 \times 5 = 16{,}777{,}216 \times 5 = 83{,}886{,}080 \approx 83.9 \text{ triệu kiến trúc.}
    \label{eq:search_space_size}
\end{equation}

Không gian gần 84 triệu kiến trúc này quá lớn để duyệt toàn bộ bằng phương pháp naïve (đánh giá từng kiến trúc độc lập sẽ cần hàng nghìn năm GPU). Đây là lý do cốt lõi tại sao AutoShot cần quy trình NAS hai giai đoạn với SuperNet và Bayesian Optimization, được trình bày tiếp theo.

\subsection{Quy trình tìm kiếm NAS}
\label{subsec:autoshot_nas}

Áp dụng khung NAS đã trình bày ở Mục~\ref{sec:nas}, AutoShot cụ thể hóa ba thành phần không gian tìm kiếm, search strategy và performance estimation thành quy trình hai lớp dưới đây.

AutoShot áp dụng One-Shot NAS kết hợp Tối ưu hóa Bayesian. Có thể nhìn quy trình ở hai lớp: \textbf{lớp tìm kiếm} gồm huấn luyện SuperNet và chọn kiến trúc bằng BO; \textbf{lớp hoàn thiện mô hình} gồm retraining, chưng cất tri thức và Weight Grafting. Hình~\ref{fig:autoshot_pipeline} tóm tắt toàn bộ luồng 5 pha từ thiết kế không gian tìm kiếm đến mô hình cuối cùng.

\begin{figure}[H]
    \centering
    \resizebox{\textwidth}{!}{
    \begin{tikzpicture}[>=Latex,font=\footnotesize]
        \node[draw, rounded corners, fill=blue!12, text width=3.0cm, minimum height=1.25cm, align=center] (p1) at (0,0) {(1) Thiết kế\\không gian tìm kiếm\\$83.9$M kiến trúc};
        \node[draw, rounded corners, fill=green!12, text width=3.2cm, minimum height=1.25cm, align=center] (p2) at (4.1,0) {(2) Huấn luyện SuperNet\\Single-Path\\One-Shot};
        \node[draw, rounded corners, fill=orange!18, text width=3.2cm, minimum height=1.25cm, align=center] (p3) at (8.4,0) {(3) Bayesian Optimization\\GP + PoF\\chọn $\alpha^*$};
        \node[draw, rounded corners, fill=purple!12, text width=3.0cm, minimum height=1.25cm, align=center] (p4) at (12.4,0) {(4) Retraining\\kiến trúc $\alpha^*$};
        \node[draw, rounded corners, fill=teal!14, text width=3.2cm, minimum height=1.25cm, align=center] (p5) at (16.7,0) {(5) Tinh chỉnh cuối\\chưng cất tri thức\\+ Weight Grafting};

        \draw[->, thick] (p1) -- (p2);
        \draw[->, thick] (p2) -- (p3);
        \draw[->, thick] (p3) -- (p4);
        \draw[->, thick] (p4) -- (p5);
        \node[align=center] at (8.4,-1.4) {\footnotesize Đầu ra: mô hình AutoShot tối ưu cho bài toán SBD video ngắn};
    \end{tikzpicture}
    }
    \caption[Sơ đồ 5 pha quy trình NAS trong AutoShot]{Sơ đồ 5 pha của quy trình NAS trong AutoShot: (1) Thiết kế không gian tìm kiếm với 83.9M kiến trúc; (2) Huấn luyện SuperNet theo chiến lược Single-Path One-Shot; (3) Tìm kiếm kiến trúc tối ưu bằng Bayesian Optimization (GP + PoF); (4) Huấn luyện lại kiến trúc $\alpha^*$; (5) Tinh chỉnh bằng chưng cất tri thức và Weight Grafting.}
    \label{fig:autoshot_pipeline}
\end{figure}

\subsubsection{Giai đoạn 1: Huấn luyện SuperNet (Single-Path One-Shot)}
\label{subsubsec:supernet_training}

SuperNet là một mạng duy nhất triển khai tất cả 83.9 triệu kiến trúc trong không gian tìm kiếm thông qua cơ chế chia sẻ trọng số (weight sharing). Mỗi trong 6 vị trí khối tích chập, SuperNet duy trì đầy đủ trọng số cho tất cả 16 cấu hình khả thi; tại mỗi bước huấn luyện, chỉ một đường đi (path) được kích hoạt và cập nhật.

Chiến lược \textit{Single-Path One-Shot} hoạt động như sau: tại mỗi bước huấn luyện, lấy mẫu ngẫu nhiên một vector cấu hình $\alpha = [\alpha_1, \alpha_2, \ldots, \alpha_6, \alpha_7] \in \mathcal{A}$ theo phân phối đồng đều, trong đó $\alpha_i$ xác định loại và siêu tham số của khối thứ $i$. Chỉ trọng số của các lớp tương ứng với $\alpha$ được kích hoạt trong bước forward/backward, còn phần còn lại không được cập nhật. Về mặt toán học, cập nhật tham số tại bước $n$ là:
\begin{equation}
    \theta_{\mathrm{super}}[\alpha^{(n)}]
    \;\leftarrow\;
    \theta_{\mathrm{super}}[\alpha^{(n)}] - lr \cdot \nabla_{\theta[\alpha^{(n)}]} \mathcal{L}(\alpha^{(n)}, \mathcal{B}^{(n)}),
    \label{eq:supernet_update}
\end{equation}
trong đó $\alpha^{(n)}$ là kiến trúc được lấy mẫu tại bước $n$, $\mathcal{B}^{(n)}$ là batch dữ liệu, và $\mathcal{L}$ là hàm mất mát giống phương trình~\eqref{eq:transnetv2_loss}.

Chiến lược lấy mẫu đồng đều đảm bảo mọi kiến trúc con được huấn luyện với xác suất bằng nhau, tránh thiên lệch về một loại khối nhất định. Điều này quan trọng để trọng số chia sẻ $\theta_{\mathrm{super}}$ phản ánh khả năng chung của tất cả kiến trúc, không bị áp đảo bởi một số ít kiến trúc chiếm ưu thế. Sau khi huấn luyện 12 epoch, $\theta_{\mathrm{super}}$ đóng vai trò oracle đánh giá nhanh: cho bất kỳ kiến trúc $\alpha$, điểm F1 ước tính được tính bằng cách nạp $\theta_{\mathrm{super}}$ và suy diễn trên tập xác thực mà không cần huấn luyện lại.

Thuật toán~\ref{alg:supernet_training} mô tả quy trình huấn luyện SuperNet; điểm mới so với TransNetV2 là bước lấy mẫu kiến trúc $\alpha$ ở đầu mỗi vòng lặp (iteration) và chỉ cập nhật có chọn lọc theo $\alpha$.

Dù trọng số chia sẻ (shared) không tối ưu cho bất kỳ kiến trúc con cụ thể nào, \textit{thứ tự xếp hạng} giữa các kiến trúc con khi dùng trọng số chia sẻ tương quan cao với thứ tự khi huấn luyện độc lập~\cite{ren2020_nas_survey}, đủ để BO hướng dẫn tìm kiếm. Sau 12 epoch, hàm đánh giá nhanh chỉ cần nạp $\theta_{\mathrm{super}}$ vào $\alpha$ và suy diễn trên tập xác thực (Thuật toán~\ref{alg:transnetv2_inference}), mất vài giây thay vì nhiều giờ.

\begin{algorithm}[H]
\caption{SuperNet Training (Single-Path One-Shot)}
\label{alg:supernet_training}
\begin{algorithmic}[1]
\Function{TrainSuperNet}{$\mathcal{S},\; epochs = 12$}
\State Initialize $\theta_{\mathrm{super}}$ randomly
\State Set: $lr = 0.1$, $momentum = 0.9$, $batch\_size = 16$, $\lambda_1 = 5.0$, $\lambda_2 = 0.1$

\For{$epoch = 1 \dots epochs$}
    \For{$iter = 1 \dots num\_iterations$}
        \State $\alpha \gets \mathrm{SampleUniform}(\mathcal{S})$
        \Comment{Sample sub-architecture uniformly at random}
        \State $\mathcal{B} \gets \text{CreateBatch}(\text{SHOT} \cup \text{ClipShots},\; batch\_size)$
        \State $(\hat{y}^{(1)}, \hat{y}^{(2)}) \gets \text{Forward}(\theta_{\mathrm{super}},\; \mathcal{B},\; \alpha)$
        \State $\mathcal{L} \gets \lambda_1 \cdot \mathrm{CE}(y^{(1)}, \hat{y}^{(1)}) + \lambda_2 \cdot \mathrm{CE}(y^{(2)}, \hat{y}^{(2)})$
        \State $\theta_{\mathrm{super}}[\alpha] \gets \theta_{\mathrm{super}}[\alpha] - lr \cdot \nabla_{\theta_{\mathrm{super}}[\alpha]} \mathcal{L}$
        \Comment{Update only path $\alpha$}
    \EndFor
\EndFor
\State \Return $\theta_{\mathrm{super}}$
\EndFunction
\end{algorithmic}
\end{algorithm}

\subsubsection{Giai đoạn 2: Tìm kiếm kiến trúc bằng Tối ưu hóa Bayesian}
\label{subsubsec:bayesian_search}

Sau khi có $\theta_{\mathrm{super}}$, AutoShot tiến hành tìm kiến trúc tối ưu bằng Tối ưu hóa Bayesian (BO) với mô hình đại diện Quá trình Gauss (GP). BO phù hợp với bài toán này vì: (1) không gian kiến trúc rời rạc và không khả vi, (2) mỗi lần đánh giá tuy rẻ hơn huấn luyện đầy đủ nhưng vẫn tốn thời gian, và (3) cần tìm cực trị toàn cục trong không gian 83.9M điểm với ngân sách hữu hạn.

GP với kernel Hamming (đo khoảng cách giữa hai chuỗi cấu hình theo số vị trí khác biệt) duy trì phân phối hậu nghiệm $f(\alpha) \mid \mathcal{D}_n \sim \mathcal{GP}(\mu_n(\alpha), \sigma_n^2(\alpha))$ dựa trên các quan sát đã thu thập. Hàm thu thập \textit{Probability of Feasibility} (PoF) hướng dẫn lựa chọn ứng viên tiếp theo:
\begin{equation}
    \mathrm{PoF}(\alpha) = \Phi\!\left(\frac{\mu_n(\alpha) - f_{\mathrm{best}}}{\sigma_n(\alpha)}\right),
    \label{eq:pof}
\end{equation}
PoF ổn định hơn EI trong không gian rời rạc có nhiễu vì chỉ quan tâm đến xác suất vượt mốc tốt nhất, không bị thiên lệch bởi dự đoán $\mu_n$ quá lạc quan ở vùng phương sai cao. Tổng cộng 4800 đánh giá ($0.0057\%$ của 83.9M kiến trúc) đủ để tìm cấu hình gần tối ưu.

\begin{algorithm}[H]
\caption{Bayesian Architecture Search}
\label{alg:bayesian_search}
\begin{algorithmic}[1]
\Function{BayesianSearch}{$\theta_{\mathrm{super}},\; T = 100$}
\State $\mathcal{D} \gets \emptyset$; \; $f_{\mathrm{best}} \gets -\infty$
\State $GP \gets \text{InitializeGaussianProcess}(\text{kernel} = \text{Hamming})$

\State \Comment{Phase 1: Random Exploration (20 epochs)}
\For{$t = 1 \dots 20$}
    \For{$i = 1 \dots 48$}
        \State $\alpha \gets \mathrm{SampleRandom}(\mathcal{S})$
        \State $F1 \gets \text{QuickEvaluate}(\alpha,\; \theta_{\mathrm{super}},\; \mathcal{D}_{\mathrm{val}})$
        \State $\mathcal{D} \gets \mathcal{D} \cup \{(\alpha, F1)\}$; \; $f_{\mathrm{best}} \gets \max(f_{\mathrm{best}}, F1)$
        \State $GP.\text{Update}(\mathcal{D})$
    \EndFor
\EndFor

\State \Comment{Phase 2: Bayesian Optimization (80 epochs)}
\For{$t = 21 \dots T$}
    \For{$i = 1 \dots 48$}
        \State $(\mu, \sigma) \gets GP.\text{Predict}(\mathcal{S})$
        \Comment{Posterior distribution over search space}
        \State $\alpha_{\mathrm{next}} \gets \arg\max_{\alpha \in \mathcal{S}} \;\Phi\!\left(\dfrac{\mu(\alpha) - f_{\mathrm{best}}}{\sigma(\alpha)}\right)$
        \State $F1 \gets \text{QuickEvaluate}(\alpha_{\mathrm{next}},\; \theta_{\mathrm{super}},\; \mathcal{D}_{\mathrm{val}})$
        \State $\mathcal{D} \gets \mathcal{D} \cup \{(\alpha_{\mathrm{next}}, F1)\}$; \; $f_{\mathrm{best}} \gets \max(f_{\mathrm{best}}, F1)$
        \State $GP.\text{Update}(\mathcal{D})$
    \EndFor
\EndFor

\State $\alpha^* \gets \arg\max_{(\alpha, F1) \in \mathcal{D}} F1$
\State \Return $\alpha^*$
\EndFunction
\end{algorithmic}
\end{algorithm}

\subsubsection{Chi phí tính toán}
\label{subsec:autoshot_complexity}

Tổng thể, quy trình AutoShot đắt hơn TransNetV2 vì phải trả chi phí tìm kiếm trước khi có mô hình cuối. Bảng~\ref{tab:autoshot_complexity} chuẩn hóa các con số theo cùng dải ước lượng được dùng trong phần so sánh.

\begin{table}[H]
\centering
\caption{Chi phí tính toán ước tính của từng pha trong AutoShot}
\label{tab:autoshot_complexity}
\small
\begin{tabularx}{\textwidth}{lYYr}
\toprule
\textbf{Pha} & \textbf{Mô tả} & \textbf{Phần cứng} & \textbf{Thời gian ước tính} \\
\midrule
1. Huấn luyện SuperNet & 12 epoch, single-path & 8× V100 & $\approx$ 12--16 giờ \\
2. Tìm kiếm kiến trúc & 4800 đánh giá (GP + PoF) & 8× V100 & $\approx$ 40--50 giờ \\
3. Huấn luyện lại (Retraining) & 12 epoch với $\alpha^*$ & 1× V100 & $\approx$ 12 giờ \\
4. Chưng cất tri thức (KD) & 12 epoch, student học từ teacher & 1× V100 & $\approx$ 12 giờ \\
5. Ghép trọng số & Tính entropy + nội suy trọng số & CPU & $<$ 1 giờ \\
\midrule
\textbf{Tổng} & & & \textbf{$\approx$ 76--91 giờ} \\
\bottomrule
\end{tabularx}
\end{table}

So với TransNetV2 ($\approx17$ giờ, 1 GPU), AutoShot tốn 4--5 lần thời gian và đòi hỏi nhiều GPU ở giai đoạn tìm kiếm. Đây là đánh đổi cốt lõi của NAS: chi phí trả trước lớn, nhưng mô hình cuối cùng có chi phí suy diễn xấp xỉ TransNetV2 vì chỉ còn một kiến trúc con $\alpha^*$ được chọn ra. SuperNet cần $\approx16$ GB VRAM trong pha huấn luyện, còn mô hình cuối chỉ cần $\approx2$ GB.

\subsection{Huấn luyện lại và tinh chỉnh mô hình}
\label{subsec:autoshot_retrain}

Sau khi BO chọn được kiến trúc $\alpha^*$, AutoShot không dùng trực tiếp trọng số chia sẻ (shared) của SuperNet làm mô hình cuối. Kiến trúc này được huấn luyện lại từ đầu trong 12 epoch với cùng hàm mục tiêu như TransNetV2 (phương trình~\eqref{eq:transnetv2_loss}). Bước huấn luyện lại giúp loại bỏ nhiễu do chia sẻ trọng số (weight sharing) gây ra: trong SuperNet, mỗi đường đi chỉ được cập nhật rời rạc theo lấy mẫu; trong mô hình cuối, toàn bộ tham số của $\alpha^*$ được tối ưu liên tục cho một kiến trúc cố định.

Sau khi huấn luyện lại, AutoShot áp dụng hai kỹ thuật tinh chỉnh không làm thay đổi kiến trúc: chưng cất tri thức và Weight Grafting. Vì vậy, phần này cần được hiểu là giai đoạn \textit{hoàn thiện trọng số}, không phải tiếp tục tìm kiếm kiến trúc.

\subsubsection{chưng cất tri thức — trích xuất tri thức}
\label{subsubsec:kd}

chưng cất tri thức (KD) \cite{hinton2015distilling} sử dụng mô hình ứng viên tốt nhất sau bước huấn luyện lại làm \textit{giáo viên} (teacher), hướng dẫn quá trình học của \textit{học sinh} (student) — cũng là kiến trúc $\alpha^*$ nhưng khởi tạo ngẫu nhiên. Ý tưởng cốt lõi: thay vì học từ nhãn cứng (nhãn cứng) nhị phân $\{0, 1\}$, student học từ nhãn mềm (nhãn mềm) của teacher.

Nhãn mềm $\tilde{y}^{(j)}_{i,t}$ là xác suất dự đoán của teacher tại khung $t$ của chuỗi $i$, nhánh $j \in \{1, 2\}$. Hàm mất mát KD:
\begin{equation}
    \mathcal{L}_{\mathrm{distill}}
    = -\frac{1}{KM} \sum_{i=1}^{K} \sum_{t=1}^{M}
      \Bigl[
        \lambda_1 \cdot \tilde{y}^{(1)}_{i,t} \log \hat{y}^{(1)}_{i,t}
        + \lambda_2 \cdot \tilde{y}^{(2)}_{i,t} \log \hat{y}^{(2)}_{i,t}
      \Bigr],
    \label{eq:kd_loss}
\end{equation}
trong đó $\tilde{y}^{(j)}$ là nhãn mềm của teacher (cố định, không có gradient qua teacher), $\hat{y}^{(j)}$ là dự đoán của student, và $\lambda_1 = 5.0$, $\lambda_2 = 0.1$ như trong phương trình~\eqref{eq:transnetv2_loss}.

So sánh với phương trình~\eqref{eq:transnetv2_loss}, sự khác biệt duy nhất là nhãn ground-truth $y^{(j)}$ được thay bằng nhãn mềm $\tilde{y}^{(j)}$ từ teacher. Nhãn mềm mang nhiều thông tin hơn nhãn nhị phân: chúng phản ánh độ không chắc chắn của teacher, đặc biệt tại các khung hình thuộc vùng chuyển cảnh dần hay các chuyển cảnh đột ngột có nền màu tương đồng. Nhờ đó, student học được phân phối xác suất trơn tru hơn, ít kích hoạt cực đoan hơn và bền vững hơn trước các biến thể rìa ranh giới — đây là cốt lõi của "giảng dạy về sự không chắc chắn" trong KD.

Trong AutoShot, teacher và student có \textit{cùng kiến trúc} $\alpha^*$ — đây là trường hợp KD tự học (self-distillation). Khác với KD truyền thống (teacher lớn hơn student), tự học giúp student thoát khỏi cực trị cục bộ mà retraining trước đó có thể đã rơi vào, do điểm khởi đầu ngẫu nhiên khác nhau.

\subsubsection{Weight Grafting — ghép trọng số theo entropy}
\label{subsubsec:grafting}

Weight Grafting là kỹ thuật kết hợp trọng số từ hai mô hình dựa trên entropy của phân phối trọng số, giả định rằng lớp có entropy trọng số cao hơn đã học được biểu diễn phong phú và đa dạng hơn.

Cho hai mô hình $M_1$ (sau retraining) và $M_2$ (sau KD) với trọng số lớp $l$ lần lượt là $W_1^{(l)}$ và $W_2^{(l)}$, entropy được ước tính bằng histogram $B = 10$ bin trên phân phối của các giá trị trọng số:
\begin{equation}
    H(W^{(l)}) = -\sum_{b=1}^{B} p_b \log p_b, \quad p_b = \frac{|\{w \in W^{(l)} : w \in \mathrm{bin}_b\}|}{|W^{(l)}|},
    \label{eq:weight_entropy}
\end{equation}
với $p_b$ là xác suất ước tính của bin $b$ trong histogram. Hệ số ghép $\alpha^{(l)}$ xác định tỉ lệ đóng góp của $M_2$ tại lớp $l$ theo công thức arctan liên tục:
\begin{equation}
    \alpha^{(l)} = A \cdot \arctan\!\bigl(c \cdot (H(W_2^{(l)}) - H(W_1^{(l)}))\bigr) + 0.5,
    \label{eq:grafting_alpha}
\end{equation}
trong đó $A = 0.4$ và $c = 1.0$ là siêu tham số điều khiển độ nhạy và tốc độ chuyển tiếp. Hàm arctan đảm bảo $\alpha^{(l)} \in [0.5 - A\pi/2,\; 0.5 + A\pi/2] \approx [0.5 - 0.628,\; 0.5 + 0.628]$, giới hạn ảnh hưởng tối đa của bất kỳ mô hình nào. Trọng số ghép cuối cùng:
\begin{equation}
    W_{\mathrm{grafted}}^{(l)} = \alpha^{(l)} \cdot W_2^{(l)} + (1 - \alpha^{(l)}) \cdot W_1^{(l)}.
    \label{eq:grafting}
\end{equation}

Khi $H(W_2^{(l)}) > H(W_1^{(l)})$, hàm arctan dương, $\alpha^{(l)} > 0.5$ — ưu tiên trọng số $M_2$ tại lớp đó; ngược lại ưu tiên $M_1$. Điều này cho phép Weight Grafting thích ứng linh hoạt theo từng lớp: lớp nào $M_2$ học tốt hơn thì dùng nhiều $M_2$, và ngược lại.

Hàm arctan đảm bảo tính bão hòa: khi chênh lệch entropy rất lớn, $\alpha^{(l)}$ tiệm cận $0.5 \pm A\pi/2$ thay vì để một mô hình hoàn toàn lấn át mô hình kia. Tham số $c$ điều chỉnh độ dốc chuyển tiếp.

Một điểm thực tế quan trọng: trong AutoShot, Weight Grafting được áp dụng giữa ba mô hình $[\theta_{\mathrm{retrain}}, \theta_{\mathrm{student}}, \theta_{\mathrm{retrain}}]$ (theo thứ tự ghép), không chỉ hai mô hình. Điều này tạo ra một chuỗi ghép tuần tự, trong đó kết quả ghép của hai mô hình đầu tiên tiếp tục được ghép với mô hình thứ ba. Kết hợp Weight Grafting sau KD mang lại cải thiện thêm $\approx 0.3\%$ F1 trên tập SHOT.

\subsection{Kiến trúc tối ưu của AutoShot}
\label{subsec:autoshot_optimal}

Nhóm tác giả tiến hành tìm kiếm hai lần với hai tiêu chí đánh giá khác nhau, thu được hai kiến trúc có đặc điểm bổ sung cho nhau: AutoShot@F1 tối ưu cân bằng precision--recall, còn AutoShot@Precision ưu tiên giảm dương tính giả.

Bảng~\ref{tab:autoshot_architectures} tổng hợp kiến trúc tìm được cho hai tiêu chí:

\begin{table}[H]
\centering
\caption{Kiến trúc tối ưu của AutoShot theo hai tiêu chí tối ưu}
\label{tab:autoshot_architectures}
\small
\begin{tabularx}{\textwidth}{lYY}
\toprule
\textbf{Block/Thông số} & \textbf{AutoShot@F1} (F1=84.1\%) & \textbf{AutoShot@Precision} (P=93.9\%) \\
\midrule
Block 1 & DDCNNV2 ($4F$, $nd=4$) & DDCNNV2 ($12F$, $nd=4$) \\
Block 2--4 & DDCNNV2A ($4F$, $nd=5$) ×3 & V2($8F,4$), V2B($4$), V2C($4$) \\
Block 5 & DDCNNV2 ($12F$, $nd=5$) & DDCNNV2B ($nd=5$) \\
Block 6 & DDCNNV2 ($8F$, $nd=5$) & DDCNNV2B ($nd=4$) \\
Block 7 (Transformer) & 0 lớp attention & 0 lớp attention \\
FLOPs & 37 GMACs & 30 GMACs \\
\bottomrule
\end{tabularx}
\end{table}

Cả hai kiến trúc đều chọn 0 lớp Transformer (xác nhận tích chập giãn nở đủ để mô hình hóa phụ thuộc thời gian trong video ngắn). AutoShot@F1 thiên về DDCNNV2A (shared spatial) ở lớp giữa, trong khi AutoShot@Precision dùng nhiều DDCNNV2B (Add) ở lớp sau — biểu diễn bảo thủ hơn, ít kích hoạt cực đoan, dẫn đến giảm dương tính giả. Điều này cho thấy không gian tìm kiếm đủ đa dạng để NAS tìm kiến trúc chuyên biệt cho từng tiêu chí.

\subsection{So sánh với TransNetV2}
\label{sec:comparison}

Tiểu mục này đặt TransNetV2 và AutoShot cạnh nhau nhằm tách bạch hai lớp đóng góp: các thành phần AutoShot kế thừa trực tiếp từ TransNetV2 và các thay đổi do quy trình NAS của AutoShot gốc mang lại. Cách nhìn này quan trọng vì AutoShot không thay toàn bộ quy trình, mà giữ lại nhiều quyết định thiết kế hiệu quả của TransNetV2 rồi chỉ tối ưu những điểm có khả năng ảnh hưởng mạnh đến phân phối SHOT.

\subsubsection{Bảng tổng hợp}

Bảng~\ref{tab:autoshot_vs_transnetv2} tập trung vào các chiều kỹ thuật cốt lõi: cách thiết kế kiến trúc, không gian tìm kiếm, thành phần giữ nguyên, chi phí huấn luyện và hiệu năng chính trên SHOT.

\begin{table}[H]
\centering
\caption{So sánh tổng hợp TransNetV2 và AutoShot}
\label{tab:autoshot_vs_transnetv2}
\small
\begin{tabularx}{\textwidth}{L{3.3cm}QQ}
\toprule
\textbf{Đặc điểm} & \textbf{TransNetV2} & \textbf{AutoShot} \\
\midrule
Cách thiết kế kiến trúc & Thiết kế thủ công, cố định & Tìm kiếm tự động bằng NAS \\
Không gian kiến trúc & 1 kiến trúc cố định & 83.9M ứng viên \\
Khối tích chập chính & DDCNNV2 cố định & DDCNNV2/A/B/C được chọn theo từng vị trí \\
Thành phần giữ nguyên & Sliding window, similarity branch, two-head output & Kế thừa các thành phần này từ TransNetV2 \\
Chiến lược tìm kiếm & Không có & Single-Path One-Shot + Bayesian Optimization \\
Kỹ thuật sau tìm kiếm & Huấn luyện trực tiếp & Retraining, chưng cất tri thức và Weight Grafting \\
Chi phí huấn luyện & $\approx$17 giờ trên 1$\times$V100 & $\approx$76--91 giờ trên 8$\times$V100 \\
Quy mô tham số & $\approx$6.5M & $\approx$6.7M \\
\midrule
Hiệu năng chính trên SHOT & F1 = 0.7993 (79.9\%) & \textbf{F1 = 0.8405 (84.1\%)} \\
Precision trên SHOT & Khoảng 80\% & \textbf{93.9\%} với AutoShot@Precision \\
\bottomrule
\end{tabularx}
\end{table}

\subsubsection{Nhận xét ưu và nhược điểm tương đối}

\textbf{Ưu điểm của AutoShot so với TransNetV2.} Mức tăng từ 0.7993 lên 0.8405 F1 trên SHOT là đáng kể, nhất là khi AutoShot không lớn hơn TransNetV2 theo số tham số. Điểm chính không nằm ở việc mở rộng năng lực mô hình, mà ở việc thay đổi cách chọn kiến trúc: NAS tìm cấu hình khối tích chập phù hợp hơn với phân phối video ngắn của SHOT. AutoShot@Precision cũng cho thấy cùng một không gian kiến trúc có thể được tối ưu theo mục tiêu khác, đạt precision 93.9\% khi cần giảm dương tính giả trong hệ thống biên tập bán tự động.

\textbf{Hạn chế của AutoShot.} Đánh đổi chính là chi phí NAS và độ khó tái lập. Quy trình AutoShot cần nhiều pha hơn TransNetV2, tiêu tốn khoảng 76--91 giờ trên 8 GPU V100 so với khoảng 17 giờ trên 1 GPU V100 của TransNetV2. Ngoài ra, kiến trúc tìm được gắn chặt với mục tiêu và tập validation dùng trong tìm kiếm; khả năng tổng quát hóa sang các bộ dữ liệu khác được phân tích riêng ở Chương~\ref{ch:ket_luan}.

Tóm lại, AutoShot kế thừa phần lớn quy trình hiệu quả của TransNetV2 và chỉ thay đổi cơ chế lựa chọn mạng xương sống. Điểm đóng góp chính không phải là làm mô hình lớn hơn, mà là làm cho kiến trúc khớp hơn với phân phối video ngắn. Đây là nền tảng để AutoShotV2 ở mục tiếp theo đóng băng mạng xương sống AutoShot@F1 và tập trung cải tiến tầng phân loại cùng hậu xử lý, thay vì chạy lại toàn bộ NAS.

\section{Biến thể AutoShotV2 --- Cải tiến tầng phân loại và hậu xử lý}
\label{sec:autoshotv2}

\subsection{Động lực và chiến lược tổng thể}
\label{subsec:v2_motivation}

AutoShot@F1 đạt F1 = 0.8405 với BCE chuẩn và nhãn một điểm dương (one-hot). Tầng phân loại vẫn còn hai hạn chế: \textbf{(i)} mất cân bằng lớp nghiêm trọng (ranh giới chiếm <1\% khung hình); \textbf{(ii)} tín hiệu cứng nhắc không chịu sai lệch chú thích. AutoShotV2 áp dụng chiến lược \textbf{đóng băng mạng xương sống + huấn luyện lại tầng phân loại}: khóa toàn bộ trọng số AutoShot@F1, huấn luyện bộ phân lớp mới từ đầu trên các đặc trưng đã được lưu bộ nhớ đệm. Quyết định này được thúc đẩy bởi ba lý do thực tế: \textbf{(a)} NAS đã tối ưu mạng xương sống với chi phí lớn (~76--91 giờ trên 8$\times$V100, Bảng~\ref{tab:autoshot_complexity}), nên kết quả không nên bị mất đi qua một lần huấn luyện toàn mạng nữa; \textbf{(b)} đóng băng mạng xương sống cô lập đóng góp của tầng phân loại, làm mọi cải thiện F1 đều có thể quy về thiết kế tầng phân loại; \textbf{(c)} chỉ huấn luyện tầng phân loại trên đặc trưng đã bộ nhớ đệm giúp giảm chi phí huấn luyện lại từ hàng chục giờ xuống còn vài chục phút trên một GPU phổ thông, cho phép thử nghiệm nhanh nhiều cấu hình phân rã.

\subsection{Kiến trúc tầng phân loại (ClassificationHead)}
\label{subsec:v2_head}

Tương tự AutoShot, AutoShotV2 thiết kế một tầng phân loại với hai nhánh song song đặt sau lớp trích xuất đặc trưng 4864-D:
\begin{enumerate}
    \item \textbf{fc1}: Linear(4864 $\to$ 1024) + ReLU --- nén đặc trưng về không gian thấp hơn.
    \item \textbf{Dropout}($p = 0.5$) --- chính quy hóa để tránh quá khớp trên tập huấn luyện nhỏ.
    \item \textbf{cls\_layer1}: Linear(1024 $\to$ 1) --- nhánh chính với \textbf{nhãn một điểm dương (one-hot)}, học định vị ranh giới chính xác.
    \item \textbf{cls\_layer2}: Linear(1024 $\to$ 1) --- nhánh phụ với \textbf{nhãn nhiều điểm dương (many-hot)}, học ngữ cảnh vùng quanh ranh giới.
\end{enumerate}

Toàn bộ đặc trưng mạng xương sống được \textbf{tính trước và lưu cache} vào file \texttt{.pickle} trước khi huấn luyện, giúp mỗi epoch chỉ cần chạy qua bộ phân lớp nhỏ thay vì toàn mạng mạng xương sống 55 MB, tăng tốc đáng kể vòng lặp huấn luyện.

\begin{figure}[H]
\centering
\begin{tikzpicture}[
  scale=0.88,
  transform shape,
  box/.style={draw, rounded corners, minimum width=3.6cm, minimum height=0.65cm, align=center, font=\small},
  arr/.style={-{Stealth[length=5pt]}, thick},
  node distance=0.45cm and 1.8cm
]
\node[box, fill=blue!10] (feat) {Features (4864-D)};
\node[box, fill=green!10, below=of feat] (fc1) {Linear(4864 $\to$ 1024) + ReLU};
\node[box, fill=green!10, below=of fc1] (drop) {Dropout($p=0.5$)};

\node[box, fill=orange!20, below left=0.6cm and 1.0cm of drop] (cls1) {cls\_layer1\\Linear(1024$\to$1)};
\node[box, fill=purple!20, below right=0.6cm and 1.0cm of drop] (cls2) {cls\_layer2\\Linear(1024$\to$1)};

\node[box, fill=orange!30, below=0.55cm of cls1] (oh) {$\hat{y}_1$ (one-hot)};
\node[box, fill=purple!30, below=0.55cm of cls2] (mh) {$\hat{y}_2$};

\path ($(oh.south)!0.5!(mh.south)$) coordinate (lossmid);
\node[box, fill=red!10, minimum width=7.8cm] (loss)
  at ($(lossmid)+(0,-2.45cm)$)
  {$\mathcal{L} = \mathcal{L}_{\text{focal}}(\hat{y}_1,\,y_{\text{oh}}) + 0.3\,\mathcal{L}_{\text{focal}}(\hat{y}_2,\,y_{\text{mh}})$};

\draw[arr] (feat) -- (fc1);
\draw[arr] (fc1) -- (drop);
\draw[arr] (drop.south) -- ++(0,-0.28) -| (cls1.north);
\draw[arr] (drop.south) -- ++(0,-0.28) -| (cls2.north);
\draw[arr] (cls1) -- (oh);
\draw[arr] (cls2) -- (mh);
\draw[arr] (oh.south)  -- ++(0,-0.4) -| ([xshift=-1.35cm]loss.north);
\draw[arr] (mh.south)  -- ++(0,-0.4) -| ([xshift=1.35cm]loss.north);
\end{tikzpicture}
\caption[Kiến trúc tầng phân loại trong AutoShotV2]{Kiến trúc tầng phân loại (ClassificationHead) trong AutoShotV2. Tầng phân loại này nhận đặc trưng 4864-D từ mạng xương sống (đóng băng) và sinh hai đầu ra song song: nhánh một điểm dương (one-hot) định vị ranh giới chính xác và nhánh nhiều điểm dương (many-hot) học ngữ cảnh vùng biên.}
\label{fig:v2_head_arch}
\end{figure}
\FloatBarrier

\subsection{Focal Loss -- giải quyết mất cân bằng lớp}
\label{subsec:v2_focal}

Trong video ngắn, khung hình tại ranh giới chiếm chưa đến 1\% tổng số khung hình. Nếu xử lý đây như bài toán phân loại nhị phân trên từng khung, hàm Binary Cross-Entropy (BCE) chuẩn bị chi phối bởi hàng loạt mẫu âm dễ phân loại:
\begin{equation}
    \mathcal{L}_{\text{BCE}} = -\left[y\log(p) + (1-y)\log(1-p)\right]
\end{equation}
Đặt $p_t$ là xác suất gán cho lớp đúng ($p_t = p$ nếu $y=1$; $p_t = 1-p$ nếu $y=0$), BCE viết gọn thành $\mathcal{L}_{\text{BCE}} = -\log(p_t)$. Khi tổng hợp giá trị mất mát (loss) trên toàn bộ lô (batch), các mẫu âm dễ (mẫu âm dễ - khung không có chuyển cảnh, rất dễ phân loại) áp đảo về số lượng và kéo gradient theo hướng ``luôn dự đoán không có ranh giới''.

Khác với AutoShot gốc dùng BCE có trọng số $\lambda_1, \lambda_2$ trên nhãn one-hot (phương trình~\eqref{eq:transnetv2_loss}), AutoShotV2 thay BCE bằng \textbf{Focal Loss}~\cite{Lin2017FocalLoss} --- được thiết kế chuyên cho mất cân bằng lớp trong bài toán phát hiện --- với nhân tử động $(1-p_t)^{\gamma}$ làm giảm trực tiếp đóng góp của các mẫu âm dễ ngay trong từng mẫu, thay vì chỉ cân bằng giữa hai nhánh:
\begin{equation}
    \mathcal{L}_{\text{focal}} = -\alpha(1-p_t)^{\gamma}\log(p_t)
    \label{eq:focal_loss}
\end{equation}

Nhân tử $(1-p_t)^{\gamma}$ là điểm mấu chốt. Với một \textit{mẫu âm dễ} ($p_t \approx 0.99$), nhân tử gần bằng 0, đóng góp vào tổng giá trị mất mát gần như biến mất. Với một \textit{mẫu khó (hard example)} --- khung ranh giới chỉ dự đoán được $p_t \approx 0.3$ --- nhân tử vẫn lớn, buộc mô hình tiếp tục học. Focal Loss ưu tiên các mẫu khó: ranh giới mờ, chuyển cảnh dần, hay thay đổi ánh sáng dễ gây nhầm lẫn. Nếu $\gamma = 0$, hàm mất mát Focal quay về BCE có trọng số $\alpha$.

Tham số $\gamma \in \{1.0, 2.0\}$ và $\alpha \in \{0.5, 0.75\}$ được chọn qua \textbf{tìm kiếm lưới} (8 tổ hợp), đánh giá theo F1 trên tập xác thực. Cấu hình tốt nhất: $\boldsymbol{\gamma = 1.0}$, $\boldsymbol{\alpha = 0.5}$. Khác với BCE có trọng số (BCE có trọng số - chỉ điều chỉnh trọng số tĩnh giữa hai lớp) hoặc hàm mất mát cân bằng lớp (hàm mất mát cân bằng lớp - dựa trên số mẫu hữu hiệu), Focal Loss điều chỉnh \textit{động} theo độ khó của từng mẫu --- phù hợp với SBD nơi mức độ khó biến thiên mạnh giữa chuyển cảnh đột ngột, chuyển cảnh dần và các trường hợp ánh sáng đột biến. Trong tìm kiếm lưới, $\gamma=1.0$ thắng $\gamma=2.0$ do nhánh nhiều điểm dương (Mục~\ref{subsec:v2_manyhot}) đã cung cấp thêm tín hiệu giám sát cho vùng biên, nên hệ số làm sắc quá mạnh khiến gradient ở các khung biên bị triệt tiêu sớm.

\subsection{Nhãn nhiều điểm dương (many-hot) --- giám sát vùng biên}
\label{subsec:v2_manyhot}

Nhãn một điểm dương (one-hot) chuẩn chỉ gán giá trị 1 tại đúng một khung hình ranh giới và 0 cho tất cả còn lại. Cách biểu diễn này chính xác về định vị, nhưng quá cứng nhắc trong thực tiễn SBD vì hai lý do: \textbf{(i)} chú thích thủ công có thể lệch $\pm1$ frame so với ranh giới thực sự; \textbf{(ii)} chuyển cảnh dần kéo dài trên nhiều frame, nên các khung lân cận ranh giới cũng mang tín hiệu chuyển cảnh đáng kể.

AutoShotV2 bổ sung một nhánh giám sát phụ với \textbf{nhãn nhiều điểm dương (many-hot)}~\cite{Muller2019LabelSmoothing}: tất cả các khung trong phạm vi $\pm1$ frame quanh ranh giới đều được coi là dương:
\begin{equation}
    y_{\text{many-hot}}[t] =
    \begin{cases}
        1, & \exists\, t_b \text{ sao cho } |t - t_b| \le 1 \\
        0, & \text{ngược lại}
    \end{cases}
\end{equation}

Hàm mất mát kết hợp hai nhánh với trọng số khác nhau:
\begin{equation}
    \mathcal{L}_{\text{total}} = \mathcal{L}_{\text{focal}}(\hat{y}_1,\,y_{\text{one-hot}}) + 0.3\cdot\mathcal{L}_{\text{focal}}(\hat{y}_2,\,y_{\text{many-hot}})
    \label{eq:total_loss}
\end{equation}

Nhánh một điểm dương (one-hot) giữ vai trò chính: định vị ranh giới sắc nét tại đúng thời điểm chuyển cảnh. Nhánh nhiều điểm dương (many-hot) đóng vai trò phụ (hệ số 0.3): dạy mô hình rằng các khung lân cận ranh giới là những mẫu ``gần đúng'' và không nên bị phạt nặng khi dự đoán lệch nhẹ. Hệ số 0.3 được chọn nhỏ hơn 1 để nhánh phụ hỗ trợ mà không lấn át định vị chính xác.

Bán kính $\pm 1$ frame được chọn căn cứ vào phân phối sai số chú thích thực tế trên SHOT: phần lớn ranh giới thủ công được đánh lệch không quá một khung so với chuyển cảnh đột ngột thật, và chuyển cảnh dần thường kéo dài 3--5 khung sao cho ba khung trung tâm $\{t-1, t, t+1\}$ mang tín hiệu chuyển cảnh rõ nhất. Mở rộng sang $\pm 2$ frame trong thử nghiệm sơ bộ làm tăng nhiễu giám sát vì các khung biên xa hơn (như $t \pm 2$) thường đã thuộc shot lân cận, gây nhầm chuyển cảnh đột ngột với chuyển cảnh dần và làm giảm precision.

\begin{figure}[H]
\centering
\begin{tikzpicture}[font=\small]
\node[anchor=east] at (-0.25, 0.7) {Một điểm dương:};
\foreach \i/\c/\v in {
  0/white/0, 1/white/0, 2/white/0,
  3/orange!40/1,
  4/white/0, 5/white/0, 6/white/0}
{
  \draw[fill=\c] (\i*1.05-0.42, 0.25) rectangle (\i*1.05+0.42, 1.15);
  \node at (\i*1.05, 0.7) {\v};
}
\node[anchor=east] at (-0.25, -0.7) {Nhiều điểm dương:};
\foreach \i/\c/\v in {
  0/white/0, 1/white/0,
  2/purple!30/1, 3/purple!30/1, 4/purple!30/1,
  5/white/0, 6/white/0}
{
  \draw[fill=\c] (\i*1.05-0.42, -1.15) rectangle (\i*1.05+0.42, -0.25);
  \node at (\i*1.05, -0.7) {\v};
}
\foreach \i/\lbl in {0/{$t{-}3$},1/{$t{-}2$},2/{$t{-}1$},3/{$t$},4/{$t{+}1$},5/{$t{+}2$},6/{$t{+}3$}} {
  \node[font=\scriptsize] at (\i*1.05, -1.45) {\lbl};
}
\draw[-{Stealth[length=4.5pt]}, dashed, gray!70] (3*1.05, 1.6) -- (3*1.05, 1.2);
\node[font=\scriptsize, gray!80] at (3*1.05, 1.8) {boundary};
\end{tikzpicture}
\caption[Minh họa nhãn một điểm dương và nhiều điểm dương]{Minh họa nhãn một điểm dương (one-hot) và nhiều điểm dương (many-hot). Nhãn một điểm dương chỉ gán 1 tại đúng khung ranh giới $t$; nhãn nhiều điểm dương mở rộng thêm $t\pm1$, cung cấp tín hiệu gradient dày đặc hơn quanh vùng ranh giới.}
\label{fig:v2_onehot_manyhot}
\end{figure}
\FloatBarrier

\subsection{Gaussian Smoothing --- triệt tiêu nhiễu đơn lẻ}
\label{subsec:v2_gaussian}

Sau khi mô hình sinh chuỗi xác suất $p[t]$, tín hiệu này thường chứa các đỉnh nhọn hẹp do nhiễu cục bộ: một khung đơn lẻ có xác suất cao bất thường vì flash sáng, rung máy, hoặc biến đổi màu sắc đột ngột. Nếu ngưỡng hóa trực tiếp, các đỉnh nhiễu dễ bị nhận nhầm là ranh giới cảnh. AutoShotV2 áp dụng \textbf{Gaussian smoothing 1D} với $\sigma = 2.0$ lên chuỗi xác suất trước bước ngưỡng hóa:
\begin{equation}
    \hat{p}[t] = \sum_{k} G_\sigma(k)\cdot p[t-k], \quad G_\sigma(k) = \frac{1}{\sqrt{2\pi}\,\sigma}\,e^{-k^2/(2\sigma^2)}
    \label{eq:gaussian_smooth}
\end{equation}

Phép biến đổi này là \textit{lấy trung bình có trọng số theo trục thời gian}: xác suất tại $t$ sau làm mượt phụ thuộc vào các khung lân cận, với trọng số giảm dần theo dạng chuông Gaussian. Hạt nhân Gaussian được chuẩn hóa nên chủ yếu \textit{phân phối lại} xác suất trong vùng lân cận.

Nguyên tắc hoạt động: ranh giới thật tạo phản hồi nhất quán trên vài khung liên tiếp --- sau smoothing, các khung ``hỗ trợ lẫn nhau'' và đỉnh xác suất được giữ lại. Nhiễu đơn lẻ không được hàng xóm ủng hộ bị kéo xuống đáng kể. Tham số $\sigma = 2.0$ được chốt sau khi quét $\sigma \in \{1.0, 1.5, 2.0, 2.5, 3.0\}$ trên tập xác thực của SHOT: giá trị nhỏ hơn ($\sigma \le 1.5$) chưa đủ triệt nhiễu flash, trong khi giá trị lớn hơn ($\sigma \ge 2.5$) làm phẳng cả các đỉnh ranh giới thật và đẩy ngưỡng tối ưu xuống vùng quá nhỏ, kém ổn định. So với bộ lọc trung vị (median filter) --- một lựa chọn hậu xử lý phổ biến khác --- hạt nhân Gaussian (kernel Gaussian) giữ được dạng đỉnh khi ranh giới được nhiều khung liên tiếp hỗ trợ và tránh hiện tượng cắt phẳng đột ngột, phù hợp hơn với phân phối xác suất sigmoid của tầng phân loại.

\subsection{Hiệu chỉnh nhiệt độ (Temperature Scaling) --- hiệu chỉnh xác suất}
\label{subsec:v2_temperature}

Sau huấn luyện, phân phối xác suất đầu ra của mô hình có thể chưa được hiệu chỉnh tốt (poorly calibrated): mô hình có thể \textit{thiếu tự tin (under-confident)} (xác suất bị nén về gần 0.5 dù nhiều mẫu khá chắc chắn) hoặc \textit{quá tự tin (over-confident)}. Trong SBD, hiệu chỉnh xác suất trực tiếp ảnh hưởng đến việc chọn ngưỡng tối ưu.

AutoShotV2 áp dụng \textbf{hiệu chỉnh nhiệt độ}~\cite{Guo2017Calibration}: tìm hệ số vô hướng $T^*$ tối thiểu hóa mất mát NLL (Negative Log-Likelihood) trên tập xác thực mà không huấn luyện lại bất kỳ tham số nào:
\begin{equation}
    T^* = \arg\min_{T \in [0.1,\,10.0]} \mathcal{L}_{\text{BCE}}\!\left(\sigma\!\left(\tfrac{z}{T}\right),\,y\right)
    \label{eq:temperature_scaling}
\end{equation}
trong đó $z$ là logit trước sigmoid. Điểm quan trọng: hiệu chỉnh nhiệt độ \textit{không thay đổi thứ tự tương đối} giữa các dự đoán (phép chia cho $T > 0$ là đơn điệu), chỉ hiệu chỉnh mức độ tự tin. Nếu $T > 1$, phân phối bị làm phẳng; nếu $T < 1$, phân phối trở nên sắc hơn và các giá trị bị đẩy xa khỏi 0.5.

Giá trị tối ưu trên tập SHOT: $T^* = 0.3878 < 1$, cho thấy mô hình đang \textit{thiếu tự tin (under-confident)} --- nhiều dự đoán đúng nhưng xác suất chưa đủ ``dứt khoát''. Chia logit cho $T^* < 1$ làm sắc hóa phân phối, đẩy các dự đoán rõ ràng ra xa vùng trung tính. Hệ quả: ngưỡng tối ưu dịch từ 0.296 (AutoShot) xuống 0.02 (AutoShotV2).

Khoảng tìm kiếm $T \in [0.1, 10.0]$ được chọn theo khuyến nghị của Guo et al.\ \cite{Guo2017Calibration} cho bài toán phân loại nhị phân: cận dưới $0.1$ cho phép sắc hóa rất mạnh khi mô hình thiếu tự tin (under-confident), còn cận trên $10.0$ phủ trường hợp quá tự tin (over-confident) mà không gặp rủi ro chia cho giá trị quá lớn làm gradient triệt tiêu. So với hiệu chỉnh Platt (Platt scaling - học hai tham số affine $a, b$ trên logit) hoặc hồi quy đơn điệu (isotonic regression - hàm bậc thang không tham số), hiệu chỉnh nhiệt độ chỉ có một tham số duy nhất nên không có nguy cơ quá khớp (quá khớp) trên tập xác thực nhỏ và giữ nguyên thứ tự xếp hạng dự đoán --- một tính chất then chốt để ngưỡng hóa nhị phân của SBD vẫn ổn định sau hiệu chỉnh.

\subsection{Quy trình huấn luyện lại và thuật toán tổng thể}
\label{subsec:v2_training}

Tầng phân loại (ClassificationHead) được huấn luyện lại theo chiến lược hai pha~\cite{Loshchilov2016SGDR}: khởi động tuyến tính trong 4 epoch (tốc độ học - LR từ 0 đến $10^{-5}$) rồi giảm nhiệt cosine đến tối đa 20 epoch. Cấu hình chạy cuối không dùng cơ chế dừng sớm, nhằm giữ cùng ngân sách huấn luyện cho các biến thể thực nghiệm phân rã.

\section{Tổng kết chương}
\label{sec:chap3_conclusion}

Chương này đã trình bày hướng tiếp cận kế thừa AutoShot và cải tiến ở tầng đầu ra để hình thành AutoShotV2. Các nội dung chính gồm nền tảng NAS và TransNetV2, mô hình AutoShot gốc, kiến trúc tầng phân loại mới, Focal Loss, nhãn many-hot, Gaussian smoothing, hiệu chỉnh nhiệt độ và quy trình huấn luyện lại. Chương~\ref{ch:bo_du_lieu_va_thuc_nghiem} tiếp theo mô tả dữ liệu, thiết lập thực nghiệm và kết quả đánh giá.
